% figs.tex —— 由 ch6/make_figs.py 產生，請勿手動修改。
% 重新產生：python ch6/make_figs.py（第 6 章，動畫 W601–W607）

\newcommand{\SixFigSmall}{%
\begin{tikzpicture}[x=1.45cm,y=1.45cm,>=Stealth]
\draw[gray!25,step=1] (-0.4,-0.6) grid (2.75,1.45);
\draw[cExact,thick,rotate around={18.435:(1.2000,0.4000)}] (1.2000,0.4000) ellipse [x radius=1.5492, y radius=0.8944];
\draw[cBar,thick] (-0.2697,-0.0899) -- (2.6697,0.8899) node[pos=0.93,below right,font=\scriptsize,cBar,inner sep=1pt] {$2R_1$};
\draw[cBar,thick] (1.4828,-0.4485) -- (0.9172,1.2485) node[pos=0.06,right,font=\scriptsize,cBar,inner sep=1pt] {$2R_2$};
\draw[gray,densely dashed] (1.200,0.400) -- (0,0);
\draw[gray,densely dashed] (1.200,0.400) -- (1,0);
\draw[gray,densely dashed] (1.200,0.400) -- (2,0);
\draw[gray,densely dashed] (1.200,0.400) -- (2,1);
\draw[gray,densely dashed] (1.200,0.400) -- (1,1);
\draw[->,cWalk,very thick,shorten >=2.5pt,shorten <=2.5pt] (0, 0) -- (1, 0);
\draw[->,cWalk,very thick,shorten >=2.5pt,shorten <=2.5pt] (1, 0) -- (2, 0);
\draw[->,cWalk,very thick,shorten >=2.5pt,shorten <=2.5pt] (2, 0) -- (2, 1);
\draw[->,cWalk,very thick,shorten >=2.5pt,shorten <=2.5pt] (2, 1) -- (1, 1);
\fill (0,0) circle (1.7pt) node[below left,font=\scriptsize,inner sep=2pt] {$\tfrac{8}{5}$};
\fill (1,0) circle (1.7pt) node[below,font=\scriptsize,inner sep=2pt] {$\tfrac{1}{5}$};
\fill (2,0) circle (1.7pt) node[below right,font=\scriptsize,inner sep=2pt] {$\tfrac{4}{5}$};
\fill (2,1) circle (1.7pt) node[above right,font=\scriptsize,inner sep=2pt] {$1$};
\fill (1,1) circle (1.7pt) node[left,font=\scriptsize,inner sep=2pt] {$\tfrac{2}{5}$};
\draw[thick] (1.1300,0.3300) -- (1.2700,0.4700) (1.1300,0.4700) -- (1.2700,0.3300);
\node[font=\scriptsize,anchor=north west,inner sep=1pt] at (1.240,0.360) {重心};
\end{tikzpicture}}

\newcommand{\SixFigMain}{%
\begin{tikzpicture}[x=0.105cm,y=0.105cm,>=Stealth]
\draw[cWalk!85!black,line width=0.35pt,line join=round] (0,0) -- (-1,0) -- (0,0) -- (0,-1) -- (1,-1) -- (-1,-1) -- (0,-1) -- (0,-2) -- (0,-1) -- (1,-1) -- (1,-2) -- (-1,-2) -- (-1,-1) -- (-2,-1) -- (-1,-1) -- (-1,-2) -- (-1,-1) -- (-2,-1) -- (-2,-2) -- (-2,-1) -- (-1,-1) -- (-1,-2) -- (0,-2) -- (-2,-2) -- (-2,-3) -- (-2,-2) -- (-3,-2) -- (-3,-3) -- (-3,-2) -- (-4,-2) -- (-2,-2) -- (-2,-1) -- (-3,-1) -- (-3,0) -- (-4,0) -- (-4,1) -- (-3,1) -- (-3,-2) -- (-3,-1) -- (-2,-1) -- (-3,-1) -- (-2,-1) -- (-4,-1) -- (-3,-1) -- (-4,-1) -- (-3,-1) -- (-3,0) -- (-1,0) -- (-1,1) -- (-1,-2) -- (-2,-2) -- (-1,-2) -- (-2,-2) -- (-2,-1) -- (-2,-4) -- (-2,-1) -- (-2,-3) -- (-4,-3) -- (-4,-2) -- (-4,-3) -- (-3,-3);
\draw[cWalk!85!black,line width=0.35pt,line join=round] (-3,-3) -- (-3,-4) -- (-5,-4) -- (-5,-3) -- (-4,-3) -- (-6,-3) -- (-6,0) -- (-5,0) -- (-5,-1) -- (-5,2) -- (-4,2) -- (-5,2) -- (-4,2) -- (-4,0) -- (-5,0) -- (-5,1) -- (-6,1) -- (-6,0) -- (-7,0) -- (-7,1) -- (-8,1) -- (-7,1) -- (-7,0) -- (-6,0) -- (-6,1) -- (-7,1) -- (-7,2) -- (-7,1) -- (-6,1) -- (-6,0) -- (-7,0) -- (-7,-1) -- (-8,-1) -- (-8,0) -- (-11,0) -- (-11,1) -- (-12,1) -- (-12,0) -- (-11,0) -- (-13,0) -- (-13,-1) -- (-16,-1) -- (-16,-4) -- (-16,-2) -- (-17,-2) -- (-17,-3) -- (-16,-3) -- (-16,-2) -- (-14,-2) -- (-14,1) -- (-15,1) -- (-15,0) -- (-14,0) -- (-15,0) -- (-13,0) -- (-13,-1) -- (-13,1) -- (-12,1) -- (-12,0) -- (-11,0) -- (-12,0);
\draw[cWalk!85!black,line width=0.35pt,line join=round] (-12,0) -- (-11,0) -- (-11,2) -- (-11,1) -- (-9,1) -- (-9,2) -- (-8,2) -- (-8,1) -- (-6,1) -- (-6,-2) -- (-6,-1) -- (-6,-2) -- (-7,-2) -- (-7,-1) -- (-7,-2) -- (-8,-2) -- (-8,-3) -- (-7,-3) -- (-7,-1) -- (-5,-1) -- (-5,0) -- (-6,0) -- (-6,-1) -- (-6,0) -- (-8,0) -- (-8,1) -- (-7,1) -- (-7,2) -- (-7,1) -- (-7,2) -- (-8,2) -- (-8,3) -- (-7,3) -- (-8,3) -- (-8,4) -- (-6,4) -- (-6,5) -- (-6,4) -- (-6,5) -- (-5,5) -- (-5,4) -- (-3,4) -- (-3,3) -- (-3,4) -- (-2,4) -- (-2,5) -- (-2,4) -- (-2,5) -- (-3,5) -- (-2,5) -- (-4,5) -- (-4,6) -- (-4,5) -- (-3,5) -- (-4,5) -- (-4,4) -- (-4,6) -- (-3,6) -- (-3,4) -- (-3,5) -- (-3,4);
\draw[cWalk!85!black,line width=0.35pt,line join=round] (-3,4) -- (-4,4) -- (-3,4) -- (-6,4) -- (-5,4) -- (-5,5) -- (-4,5) -- (-4,4) -- (-4,5) -- (-5,5) -- (-5,7) -- (-5,6) -- (-6,6) -- (-6,7) -- (-6,6) -- (-8,6) -- (-8,7) -- (-8,6) -- (-8,7) -- (-8,5) -- (-8,6) -- (-9,6) -- (-6,6) -- (-6,5) -- (-6,6) -- (-9,6) -- (-8,6) -- (-9,6) -- (-9,7) -- (-10,7) -- (-8,7) -- (-8,6) -- (-9,6) -- (-9,5) -- (-10,5) -- (-9,5) -- (-9,6) -- (-9,5) -- (-7,5) -- (-8,5) -- (-8,6) -- (-6,6) -- (-6,5) -- (0,5) -- (-1,5) -- (-1,4) -- (-2,4) -- (-2,2) -- (-2,3) -- (-3,3) -- (-3,2) -- (-3,4) -- (1,4) -- (1,5) -- (4,5) -- (3,5) -- (3,4) -- (0,4) -- (0,2) -- (0,3) -- (1,3);
\draw[cWalk!85!black,line width=0.35pt,line join=round] (1,3) -- (1,2) -- (0,2) -- (0,1) -- (-1,1) -- (0,1) -- (-1,1) -- (0,1) -- (0,2) -- (1,2) -- (0,2) -- (0,1) -- (-1,1) -- (0,1) -- (0,2) -- (-2,2) -- (-1,2) -- (-1,3) -- (-2,3) -- (-1,3) -- (-1,2) -- (0,2) -- (0,0) -- (-1,0) -- (-1,-1) -- (-1,0) -- (-2,0) -- (0,0) -- (0,-1) -- (-1,-1) -- (-1,-3) -- (-2,-3) -- (-2,-2) -- (-2,-3) -- (-1,-3) -- (-1,-4) -- (-1,0) -- (-2,0) -- (0,0) -- (0,1) -- (-1,1) -- (-1,2) -- (-3,2) -- (-2,2) -- (-2,3) -- (-2,-2) -- (-1,-2) -- (-2,-2) -- (-2,0) -- (-2,-1) -- (0,-1) -- (-3,-1) -- (1,-1) -- (-1,-1) -- (0,-1) -- (0,-3) -- (-3,-3) -- (-3,-4) -- (-3,-3) -- (-4,-3) -- (-4,-2);
\draw[cWalk!85!black,line width=0.35pt,line join=round] (-4,-2) -- (-4,-3) -- (-5,-3) -- (-5,-1) -- (-4,-1) -- (-5,-1) -- (-4,-1) -- (-5,-1) -- (-4,-1) -- (-4,-2) -- (-7,-2) -- (-7,-3) -- (-8,-3) -- (-7,-3) -- (-7,-4) -- (-8,-4) -- (-6,-4) -- (-7,-4) -- (-7,-6) -- (-7,-5) -- (-7,-6) -- (-9,-6) -- (-9,-7) -- (-9,-6) -- (-9,-8) -- (-10,-8) -- (-10,-9) -- (-7,-9) -- (-10,-9) -- (-10,-8) -- (-10,-10) -- (-10,-9) -- (-11,-9) -- (-11,-8) -- (-13,-8) -- (-13,-9) -- (-11,-9) -- (-11,-11) -- (-12,-11) -- (-12,-10) -- (-15,-10) -- (-14,-10) -- (-15,-10) -- (-15,-8) -- (-16,-8) -- (-16,-9) -- (-17,-9) -- (-17,-10) -- (-16,-10) -- (-16,-11) -- (-16,-9) -- (-16,-10) -- (-14,-10) -- (-14,-11) -- (-12,-11) -- (-12,-12) -- (-11,-12) -- (-11,-11) -- (-11,-12) -- (-11,-11) -- (-11,-13);
\draw[cWalk!85!black,line width=0.35pt,line join=round] (-11,-13) -- (-11,-12) -- (-11,-13) -- (-10,-13) -- (-10,-14) -- (-8,-14) -- (-8,-16) -- (-9,-16) -- (-9,-15) -- (-10,-15) -- (-10,-14) -- (-9,-14) -- (-9,-16) -- (-9,-15) -- (-9,-16) -- (-12,-16) -- (-12,-17) -- (-12,-16) -- (-12,-18) -- (-12,-17) -- (-13,-17) -- (-11,-17) -- (-11,-18) -- (-9,-18) -- (-9,-21) -- (-9,-20) -- (-9,-21) -- (-10,-21) -- (-10,-22) -- (-11,-22) -- (-11,-20) -- (-11,-21) -- (-11,-20) -- (-10,-20) -- (-12,-20) -- (-12,-19) -- (-12,-21) -- (-12,-20) -- (-13,-20) -- (-13,-19) -- (-14,-19) -- (-14,-18) -- (-16,-18) -- (-16,-17) -- (-16,-18) -- (-16,-17) -- (-16,-19) -- (-15,-19) -- (-15,-18) -- (-16,-18) -- (-16,-19) -- (-16,-18) -- (-16,-20) -- (-16,-19) -- (-16,-21) -- (-16,-19) -- (-18,-19) -- (-18,-20) -- (-18,-19) -- (-19,-19) -- (-19,-20);
\draw[cWalk!85!black,line width=0.35pt,line join=round] (-19,-20) -- (-20,-20) -- (-20,-21) -- (-19,-21) -- (-20,-21) -- (-20,-22) -- (-20,-21) -- (-21,-21) -- (-21,-19) -- (-23,-19) -- (-22,-19) -- (-23,-19) -- (-22,-19) -- (-23,-19) -- (-23,-18) -- (-22,-18) -- (-22,-17) -- (-23,-17) -- (-23,-18) -- (-21,-18) -- (-21,-17) -- (-20,-17) -- (-21,-17) -- (-20,-17) -- (-21,-17) -- (-21,-16) -- (-21,-18) -- (-22,-18) -- (-22,-16) -- (-23,-16) -- (-22,-16) -- (-22,-15) -- (-22,-16) -- (-23,-16) -- (-23,-17) -- (-19,-17) -- (-19,-18) -- (-19,-17) -- (-19,-19) -- (-19,-17) -- (-19,-18) -- (-20,-18) -- (-20,-20) -- (-20,-19) -- (-19,-19) -- (-20,-19) -- (-20,-17) -- (-20,-18) -- (-19,-18) -- (-20,-18) -- (-20,-19) -- (-20,-18) -- (-21,-18) -- (-21,-19) -- (-23,-19) -- (-23,-18) -- (-22,-18) -- (-22,-17) -- (-20,-17) -- (-20,-18) -- (-20,-17);
\draw[cWalk!85!black,line width=0.35pt,line join=round] (-20,-17) -- (-22,-17) -- (-22,-16) -- (-21,-16) -- (-23,-16) -- (-23,-15) -- (-25,-15) -- (-24,-15) -- (-25,-15) -- (-24,-15) -- (-26,-15) -- (-25,-15) -- (-27,-15) -- (-27,-16) -- (-27,-14) -- (-27,-15) -- (-28,-15) -- (-28,-16) -- (-29,-16) -- (-29,-18) -- (-26,-18) -- (-26,-19) -- (-25,-19) -- (-25,-20) -- (-25,-19) -- (-25,-20) -- (-24,-20) -- (-24,-19) -- (-25,-19) -- (-25,-18) -- (-23,-18) -- (-25,-18) -- (-24,-18) -- (-25,-18) -- (-25,-17) -- (-25,-18) -- (-26,-18) -- (-25,-18) -- (-26,-18) -- (-24,-18) -- (-24,-17) -- (-25,-17) -- (-25,-16) -- (-24,-16) -- (-24,-15) -- (-24,-16) -- (-25,-16) -- (-24,-16) -- (-24,-15) -- (-24,-16) -- (-24,-15) -- (-25,-15) -- (-25,-16) -- (-26,-16) -- (-26,-15) -- (-26,-18) -- (-27,-18) -- (-27,-15) -- (-27,-16) -- (-26,-16) -- (-27,-16);
\draw[cWalk!85!black,line width=0.35pt,line join=round] (-27,-16) -- (-27,-15) -- (-29,-15) -- (-28,-15) -- (-28,-14) -- (-28,-15) -- (-29,-15) -- (-28,-15) -- (-28,-14) -- (-28,-16) -- (-29,-16) -- (-29,-15) -- (-32,-15) -- (-31,-15) -- (-31,-16) -- (-31,-15) -- (-31,-17) -- (-31,-16) -- (-33,-16) -- (-33,-18) -- (-36,-18) -- (-35,-18) -- (-36,-18) -- (-36,-17) -- (-36,-18) -- (-36,-16) -- (-37,-16) -- (-37,-18) -- (-38,-18) -- (-36,-18) -- (-36,-20) -- (-37,-20) -- (-37,-21) -- (-38,-21) -- (-38,-20) -- (-36,-20) -- (-36,-21) -- (-35,-21) -- (-35,-20) -- (-37,-20) -- (-37,-18) -- (-38,-18) -- (-37,-18) -- (-37,-19) -- (-37,-18) -- (-37,-19) -- (-36,-19) -- (-38,-19) -- (-36,-19) -- (-37,-19) -- (-36,-19) -- (-36,-17) -- (-36,-19) -- (-35,-19) -- (-35,-20) -- (-36,-20) -- (-36,-21) -- (-34,-21) -- (-34,-22) -- (-35,-22) -- (-35,-24);
\draw[cWalk!85!black,line width=0.35pt,line join=round] (-35,-24) -- (-34,-24) -- (-34,-23) -- (-34,-24) -- (-35,-24) -- (-34,-24) -- (-35,-24) -- (-35,-23) -- (-34,-23) -- (-34,-20) -- (-35,-20) -- (-34,-20) -- (-35,-20) -- (-35,-19) -- (-35,-20) -- (-35,-18) -- (-34,-18) -- (-34,-17) -- (-32,-17) -- (-32,-18) -- (-35,-18) -- (-35,-19) -- (-35,-18) -- (-33,-18) -- (-33,-19) -- (-32,-19) -- (-34,-19) -- (-33,-19) -- (-34,-19) -- (-34,-18) -- (-36,-18) -- (-34,-18) -- (-34,-16) -- (-32,-16) -- (-32,-15) -- (-33,-15) -- (-32,-15) -- (-32,-14) -- (-33,-14) -- (-32,-14) -- (-33,-14) -- (-32,-14) -- (-32,-15) -- (-31,-15) -- (-33,-15) -- (-32,-15) -- (-32,-16) -- (-35,-16) -- (-34,-16) -- (-36,-16) -- (-36,-18) -- (-35,-18) -- (-35,-19) -- (-35,-18) -- (-35,-19) -- (-35,-18) -- (-34,-18) -- (-36,-18) -- (-36,-17) -- (-36,-18) -- (-35,-18);
\draw[cWalk!85!black,line width=0.35pt,line join=round] (-35,-18) -- (-36,-18) -- (-36,-17) -- (-37,-17) -- (-36,-17) -- (-36,-16) -- (-35,-16) -- (-36,-16) -- (-35,-16) -- (-35,-15) -- (-36,-15) -- (-36,-14) -- (-36,-15) -- (-36,-14) -- (-36,-15) -- (-37,-15) -- (-36,-15) -- (-36,-16) -- (-38,-16) -- (-37,-16) -- (-37,-17) -- (-36,-17) -- (-37,-17) -- (-37,-18) -- (-38,-18) -- (-38,-17) -- (-37,-17) -- (-37,-19) -- (-38,-19) -- (-38,-18) -- (-40,-18) -- (-40,-19) -- (-38,-19) -- (-39,-19) -- (-38,-19) -- (-38,-18) -- (-38,-19) -- (-37,-19) -- (-38,-19) -- (-38,-21) -- (-38,-20) -- (-37,-20) -- (-37,-19) -- (-35,-19) -- (-35,-17) -- (-40,-17) -- (-39,-17) -- (-39,-18) -- (-40,-18) -- (-40,-19) -- (-39,-19) -- (-39,-18) -- (-40,-18) -- (-40,-17) -- (-41,-17) -- (-41,-16) -- (-39,-16) -- (-40,-16) -- (-40,-17) -- (-40,-16) -- (-41,-16);
\draw[cWalk!85!black,line width=0.35pt,line join=round] (-41,-16) -- (-41,-17) -- (-43,-17) -- (-43,-16) -- (-44,-16) -- (-43,-16) -- (-45,-16) -- (-45,-15) -- (-45,-16) -- (-45,-15) -- (-44,-15) -- (-44,-16) -- (-45,-16) -- (-45,-12) -- (-44,-12) -- (-45,-12) -- (-44,-12) -- (-45,-12) -- (-43,-12) -- (-45,-12) -- (-45,-10) -- (-44,-10) -- (-46,-10) -- (-46,-9) -- (-45,-9) -- (-45,-8) -- (-47,-8) -- (-46,-8) -- (-46,-5) -- (-47,-5) -- (-46,-5) -- (-47,-5) -- (-47,-4) -- (-49,-4) -- (-48,-4) -- (-48,-5) -- (-49,-5) -- (-49,-7) -- (-48,-7) -- (-48,-6) -- (-49,-6) -- (-49,-4) -- (-48,-4);
\draw[cExact,thick,rotate around={28.368:(-17.8731,-8.9101)}] (-17.8731,-8.9101) ellipse [x radius=32.1728, y radius=11.2494];
\draw[cBar,thick] (-46.1825,-24.1964) -- (10.4363,6.3762) node[pos=0.93,below right,font=\scriptsize,cBar,inner sep=1pt] {$2R_1$};
\draw[cBar,thick] (-12.5282,-18.8086) -- (-23.2181,0.9885) node[pos=0.93,left,font=\scriptsize,cBar,inner sep=1pt] {$2R_2$};
\draw[thick,fill=white] (0,0) circle (1.8pt);
\draw[very thick] (-18.9731,-10.0101) -- (-16.7731,-7.8101) (-18.9731,-7.8101) -- (-16.7731,-10.0101);
\draw[|-|,gray] (-48,-29) -- (-38,-29) node[midway,below,font=\scriptsize] {10};
\end{tikzpicture}}

\newcommand{\SixTabMain}{%
\begin{tabular}{@{}lccc@{}}
\toprule
 & 東東北西 & 主例（種子 643） & 主例轉 $30^\circ$\\
\midrule
點數 $N+1$ & 5 & 1001 & 1001\\
重心 & $(\frac65,\frac25)$ & $(-17.9, -8.9)$ & $(-17.9, -8.9)$\\
$T_{11},\ T_{22},\ T_{12}$ & $\frac{14}{25},\ \frac{6}{25},\ \frac{3}{25}$ & 207.5, 82.9, 95.0 & 94.1, 196.3, 101.4\\
$R_1^2,\ R_2^2$ & $\frac35,\ \frac15$ & 258.8, 31.6 & 258.8, 31.6\\
$R_1,\ R_2$ & 0.7746, 0.4472 & 16.09, 5.62 & 16.09, 5.62\\
主軸角 & $18.43^\circ$ & $28.4^\circ$ & $58.4^\circ$\\
非球度 $\big(\frac{R_1^2-R_2^2}{R_1^2+R_2^2}\big)^2$ & $\frac14$ & 0.6117 & 0.6117\\
\bottomrule
\end{tabular}}

\newcommand{\SixTabA}{%
\begin{tabular}{@{}lccccc@{}}
\toprule
$d$ & 2 & 3 & 4 & 10 & $\infty$\\
\midrule
$A_d=\frac{4+2d}{4+5d}$ & $\frac{4}{7}$ & $\frac{10}{19}$ & $\frac{1}{2}$ & $\frac{4}{9}$ & $\frac25$\\
\quad 小數 & 0.5714 & 0.5263 & 0.5000 & 0.4444 & 0.4000\\
$A'_d$（個別值的平均） & 0.396400 & 0.394274 & 0.393829 & --- & $\to\frac25$\\
\bottomrule
\end{tabular}}

\newcommand{\SixFigHist}{%
\begin{tikzpicture}[x=9cm,y=0.0016cm,>=Stealth]
\draw[cBar,fill=cBar!25] (0.00,0) rectangle (0.05,119);
\draw[cBar,fill=cBar!25] (0.05,0) rectangle (0.10,616);
\draw[cBar,fill=cBar!25] (0.10,0) rectangle (0.15,1183);
\draw[cBar,fill=cBar!25] (0.15,0) rectangle (0.20,1631);
\draw[cBar,fill=cBar!25] (0.20,0) rectangle (0.25,1890);
\draw[cBar,fill=cBar!25] (0.25,0) rectangle (0.30,1795);
\draw[cBar,fill=cBar!25] (0.30,0) rectangle (0.35,1808);
\draw[cBar,fill=cBar!25] (0.35,0) rectangle (0.40,1621);
\draw[cBar,fill=cBar!25] (0.40,0) rectangle (0.45,1669);
\draw[cBar,fill=cBar!25] (0.45,0) rectangle (0.50,1529);
\draw[cBar,fill=cBar!25] (0.50,0) rectangle (0.55,1505);
\draw[cBar,fill=cBar!25] (0.55,0) rectangle (0.60,1305);
\draw[cBar,fill=cBar!25] (0.60,0) rectangle (0.65,1065);
\draw[cBar,fill=cBar!25] (0.65,0) rectangle (0.70,1000);
\draw[cBar,fill=cBar!25] (0.70,0) rectangle (0.75,655);
\draw[cBar,fill=cBar!25] (0.75,0) rectangle (0.80,389);
\draw[cBar,fill=cBar!25] (0.80,0) rectangle (0.85,162);
\draw[cBar,fill=cBar!25] (0.85,0) rectangle (0.90,56);
\draw[cBar,fill=cBar!25] (0.90,0) rectangle (0.95,2);
\draw[->] (0,0) -- (1.06,0) node[right,font=\scriptsize] {$\hat A$};
\draw[->] (0,0) -- (0,2150) node[above,font=\scriptsize] {條數};
\foreach \x in {0,0.2,0.4,0.6,0.8,1} {\draw (\x,0) -- (\x,-40); \node[below,font=\scriptsize] at (\x,-30) {$\x$};}
\foreach \y in {500,1000,1500} {\draw (0,\y) -- (-0.01,\y); \node[left,font=\scriptsize] at (-0.01,\y) {\y};}
\draw[cExact,very thick,dashed] (0.5231,0) -- (0.5231,2050) node[anchor=south west,font=\scriptsize,inner sep=1pt] {平均的比 0.5231};
\draw[cWalk,very thick,densely dotted] (0.3939,0) -- (0.3939,2050) node[anchor=south east,font=\scriptsize,inner sep=1pt] {個別值的平均 0.3939};
\end{tikzpicture}}

\newcommand{\SixTabStair}{%
\begin{tabular}{@{}lcccc@{}}
\toprule
特徵值 & $k=1$ & $k=2$ & $k=3$ & $k=4$\\
\midrule
樓梯 $N=3$：精確 & 0.4268 & 0.1250 & 0.0732 & \\
\quad 大 $N$ 公式 $\frac{N+1}{\pi^2k^2}$ & 0.4053 & 0.1013 & 0.0450 & \\
Kramers $N=4$：精確 & 0.5236 & 0.1447 & 0.0764 & 0.0553\\
\quad 大 $N$ 公式 & 0.5066 & 0.1267 & 0.0563 & 0.0317\\
\bottomrule
\end{tabular}
\\[6pt]
\begin{tabular}[t]{@{}lcccc@{}}
\toprule
樓梯 & $N=3$ & 10 & 100 & $\infty$\\
\midrule
$\frac{\mathrm{Tr}\,\mathbf T^2}{(\mathrm{Tr}\,\mathbf T)^2}$ & 0.5200 & 0.4150 & 0.4002 & 0.4000\\
\bottomrule
\end{tabular}
\qquad
\begin{tabular}[t]{@{}ll@{}}
\toprule
三維的比 & $R_1^2:R_2^2:R_3^2$\\
\midrule
(6.53)，$d\to\infty$ & $9:2.25:1$\\
模擬 $\langle R_1^2\rangle:\langle R_2^2\rangle:\langle R_3^2\rangle$ & $11.95:2.69:1$\\
模擬 $\langle R_1^2/R_3^2\rangle:\langle R_2^2/R_3^2\rangle$ & $13.69:2.98:1$\\
書上 (6.54) & {\color{cWrong}$13.76:3.03:1$}\\
\bottomrule
\end{tabular}}

\newcommand{\SixFigOneDee}{%
\begin{tikzpicture}[>=Stealth]
\draw[cBar,thick,fill=cBar,fill opacity=0.12] (0.000,0) -- (0.000,0.000) -- (0.190,0.000) -- (0.190,0.008) -- (0.380,0.008) -- (0.380,0.032) -- (0.570,0.032) -- (0.570,0.040) -- (0.760,0.040) -- (0.760,0.048) -- (0.950,0.048) -- (0.950,0.279) -- (1.140,0.279) -- (1.140,0.479) -- (1.330,0.479) -- (1.330,1.069) -- (1.520,1.069) -- (1.520,1.404) -- (1.710,1.404) -- (1.710,1.691) -- (1.900,1.691) -- (1.900,1.906) -- (2.090,1.906) -- (2.090,1.866) -- (2.280,1.866) -- (2.280,1.970) -- (2.470,1.970) -- (2.470,1.531) -- (2.660,1.531) -- (2.660,1.228) -- (2.850,1.228) -- (2.850,0.965) -- (3.040,0.965) -- (3.040,0.558) -- (3.230,0.558) -- (3.230,0.375) -- (3.420,0.375) -- (3.420,0.215) -- (3.610,0.215) -- (3.610,0.136) -- (3.800,0.136) -- (3.800,0.104) -- (3.990,0.104) -- (3.990,0.024) -- (4.180,0.024) -- (4.180,0.008) -- (4.370,0.008) -- (4.370,0.016) -- (4.560,0.016) -- (4.560,0.000) -- (4.750,0.000) -- (4.750,0);
\draw[cWalk,thick,fill=cWalk,fill opacity=0.12] (0.000,0) -- (0.000,0.000) -- (0.190,0.000) -- (0.190,0.000) -- (0.380,0.000) -- (0.380,0.016) -- (0.570,0.016) -- (0.570,0.056) -- (0.760,0.056) -- (0.760,0.128) -- (0.950,0.128) -- (0.950,0.319) -- (1.140,0.319) -- (1.140,0.630) -- (1.330,0.630) -- (1.330,0.989) -- (1.520,0.989) -- (1.520,1.563) -- (1.710,1.563) -- (1.710,1.795) -- (1.900,1.795) -- (1.900,1.930) -- (2.090,1.930) -- (2.090,2.042) -- (2.280,2.042) -- (2.280,1.819) -- (2.470,1.819) -- (2.470,1.364) -- (2.660,1.364) -- (2.660,1.141) -- (2.850,1.141) -- (2.850,0.830) -- (3.040,0.830) -- (3.040,0.510) -- (3.230,0.510) -- (3.230,0.359) -- (3.420,0.359) -- (3.420,0.263) -- (3.610,0.263) -- (3.610,0.144) -- (3.800,0.144) -- (3.800,0.040) -- (3.990,0.040) -- (3.990,0.000) -- (4.180,0.000) -- (4.180,0.008) -- (4.370,0.008) -- (4.370,0.008) -- (4.560,0.008) -- (4.560,0.000) -- (4.750,0.000) -- (4.750,0);
\draw[cExact,thick,fill=cExact,fill opacity=0.12] (0.000,0) -- (0.000,0.000) -- (0.190,0.000) -- (0.190,0.000) -- (0.380,0.000) -- (0.380,0.008) -- (0.570,0.008) -- (0.570,0.048) -- (0.760,0.048) -- (0.760,0.128) -- (0.950,0.128) -- (0.950,0.263) -- (1.140,0.263) -- (1.140,0.646) -- (1.330,0.646) -- (1.330,0.853) -- (1.520,0.853) -- (1.520,1.460) -- (1.710,1.460) -- (1.710,1.627) -- (1.900,1.627) -- (1.900,1.858) -- (2.090,1.858) -- (2.090,2.130) -- (2.280,2.130) -- (2.280,1.914) -- (2.470,1.914) -- (2.470,1.747) -- (2.660,1.747) -- (2.660,1.188) -- (2.850,1.188) -- (2.850,0.742) -- (3.040,0.742) -- (3.040,0.582) -- (3.230,0.582) -- (3.230,0.295) -- (3.420,0.295) -- (3.420,0.239) -- (3.610,0.239) -- (3.610,0.112) -- (3.800,0.112) -- (3.800,0.080) -- (3.990,0.080) -- (3.990,0.008) -- (4.180,0.008) -- (4.180,0.008) -- (4.370,0.008) -- (4.370,0.016) -- (4.560,0.016) -- (4.560,0.000) -- (4.750,0.000) -- (4.750,0);
\draw[->] (0.000,0) -- (4.900,0);
\draw (0.000,0) -- (0.000,2.300);
\draw (0.432,0) -- (0.432,-0.08) node[below,font=\scriptsize] {$0.6$};
\draw (1.295,0) -- (1.295,-0.08) node[below,font=\scriptsize] {$0.8$};
\draw (2.159,0) -- (2.159,-0.08) node[below,font=\scriptsize] {$1.0$};
\draw (3.023,0) -- (3.023,-0.08) node[below,font=\scriptsize] {$1.2$};
\draw (3.886,0) -- (3.886,-0.08) node[below,font=\scriptsize] {$1.4$};
\draw[gray,thick,dashed] (2.159,0) -- (2.159,2.300);
\draw[black,densely dotted,thick] (2.191,0) -- (2.191,2.300);
\node[anchor=north east,font=\scriptsize,align=right] at (4.750,2.300) {$d=100$\\2000 條};
\draw[cBar,thick,fill=cBar,fill opacity=0.12] (5.500,0) -- (5.500,0.000) -- (5.690,0.000) -- (5.690,0.000) -- (5.880,0.000) -- (5.880,0.006) -- (6.070,0.006) -- (6.070,0.032) -- (6.260,0.032) -- (6.260,0.207) -- (6.450,0.207) -- (6.450,0.536) -- (6.640,0.536) -- (6.640,0.871) -- (6.830,0.871) -- (6.830,1.246) -- (7.020,1.246) -- (7.020,1.859) -- (7.210,1.859) -- (7.210,1.620) -- (7.400,1.620) -- (7.400,1.426) -- (7.590,1.426) -- (7.590,1.400) -- (7.780,1.400) -- (7.780,1.155) -- (7.970,1.155) -- (7.970,0.813) -- (8.160,0.813) -- (8.160,0.620) -- (8.350,0.620) -- (8.350,0.387) -- (8.540,0.387) -- (8.540,0.297) -- (8.730,0.297) -- (8.730,0.155) -- (8.920,0.155) -- (8.920,0.123) -- (9.110,0.123) -- (9.110,0.065) -- (9.300,0.065) -- (9.300,0.052) -- (9.490,0.052) -- (9.490,0.006) -- (9.680,0.006) -- (9.680,0.006) -- (9.870,0.006) -- (9.870,0.013) -- (10.060,0.013) -- (10.060,0.006) -- (10.250,0.006) -- (10.250,0);
\draw[cWalk,thick,fill=cWalk,fill opacity=0.12] (5.500,0) -- (5.500,0.000) -- (5.690,0.000) -- (5.690,0.000) -- (5.880,0.000) -- (5.880,0.000) -- (6.070,0.000) -- (6.070,0.000) -- (6.260,0.000) -- (6.260,0.071) -- (6.450,0.071) -- (6.450,0.445) -- (6.640,0.445) -- (6.640,0.845) -- (6.830,0.845) -- (6.830,1.446) -- (7.020,1.446) -- (7.020,1.678) -- (7.210,1.678) -- (7.210,1.833) -- (7.400,1.833) -- (7.400,1.620) -- (7.590,1.620) -- (7.590,1.478) -- (7.780,1.478) -- (7.780,1.097) -- (7.970,1.097) -- (7.970,0.852) -- (8.160,0.852) -- (8.160,0.542) -- (8.350,0.542) -- (8.350,0.471) -- (8.540,0.471) -- (8.540,0.194) -- (8.730,0.194) -- (8.730,0.148) -- (8.920,0.148) -- (8.920,0.065) -- (9.110,0.065) -- (9.110,0.039) -- (9.300,0.039) -- (9.300,0.065) -- (9.490,0.065) -- (9.490,0.006) -- (9.680,0.006) -- (9.680,0.013) -- (9.870,0.013) -- (9.870,0.000) -- (10.060,0.000) -- (10.060,0.000) -- (10.250,0.000) -- (10.250,0);
\draw[cExact,thick,fill=cExact,fill opacity=0.12] (5.500,0) -- (5.500,0.000) -- (5.690,0.000) -- (5.690,0.000) -- (5.880,0.000) -- (5.880,0.000) -- (6.070,0.000) -- (6.070,0.000) -- (6.260,0.000) -- (6.260,0.013) -- (6.450,0.013) -- (6.450,0.181) -- (6.640,0.181) -- (6.640,0.626) -- (6.830,0.626) -- (6.830,1.471) -- (7.020,1.471) -- (7.020,1.975) -- (7.210,1.975) -- (7.210,2.130) -- (7.400,2.130) -- (7.400,1.923) -- (7.590,1.923) -- (7.590,1.594) -- (7.780,1.594) -- (7.780,1.058) -- (7.970,1.058) -- (7.970,0.703) -- (8.160,0.703) -- (8.160,0.523) -- (8.350,0.523) -- (8.350,0.336) -- (8.540,0.336) -- (8.540,0.174) -- (8.730,0.174) -- (8.730,0.065) -- (8.920,0.065) -- (8.920,0.077) -- (9.110,0.077) -- (9.110,0.006) -- (9.300,0.006) -- (9.300,0.019) -- (9.490,0.019) -- (9.490,0.013) -- (9.680,0.013) -- (9.680,0.006) -- (9.870,0.006) -- (9.870,0.000) -- (10.060,0.000) -- (10.060,0.006) -- (10.250,0.006) -- (10.250,0);
\draw[->] (5.500,0) -- (10.400,0);
\draw (5.500,0) -- (5.500,2.300);
\draw (5.500,0) -- (5.500,-0.08) node[below,font=\scriptsize] {$0$};
\draw (6.450,0) -- (6.450,-0.08) node[below,font=\scriptsize] {$0.5$};
\draw (7.400,0) -- (7.400,-0.08) node[below,font=\scriptsize] {$1$};
\draw (8.350,0) -- (8.350,-0.08) node[below,font=\scriptsize] {$1.5$};
\draw (9.300,0) -- (9.300,-0.08) node[below,font=\scriptsize] {$2$};
\draw (10.250,0) -- (10.250,-0.08) node[below,font=\scriptsize] {$2.5$};
\draw[gray,thick,dashed] (7.400,0) -- (7.400,2.300);
\draw[black,densely dotted,thick] (7.471,0) -- (7.471,2.300);
\node[anchor=north east,font=\scriptsize,align=right] at (10.250,2.300) {$d=20$\\2000 條};
\draw[cBar,thick,fill=cBar,fill opacity=0.12] (11.000,0) -- (11.000,0.002) -- (11.190,0.002) -- (11.190,0.265) -- (11.380,0.265) -- (11.380,0.897) -- (11.570,0.897) -- (11.570,1.167) -- (11.760,1.167) -- (11.760,1.219) -- (11.950,1.219) -- (11.950,1.152) -- (12.140,1.152) -- (12.140,0.981) -- (12.330,0.981) -- (12.330,0.864) -- (12.520,0.864) -- (12.520,0.737) -- (12.710,0.737) -- (12.710,0.607) -- (12.900,0.607) -- (12.900,0.508) -- (13.090,0.508) -- (13.090,0.389) -- (13.280,0.389) -- (13.280,0.363) -- (13.470,0.363) -- (13.470,0.259) -- (13.660,0.259) -- (13.660,0.251) -- (13.850,0.251) -- (13.850,0.183) -- (14.040,0.183) -- (14.040,0.134) -- (14.230,0.134) -- (14.230,0.125) -- (14.420,0.125) -- (14.420,0.089) -- (14.610,0.089) -- (14.610,0.083) -- (14.800,0.083) -- (14.800,0.064) -- (14.990,0.064) -- (14.990,0.046) -- (15.180,0.046) -- (15.180,0.047) -- (15.370,0.047) -- (15.370,0.037) -- (15.560,0.037) -- (15.560,0.027) -- (15.750,0.027) -- (15.750,0);
\draw[cWalk,thick,fill=cWalk,fill opacity=0.12] (11.000,0) -- (11.000,0.000) -- (11.190,0.000) -- (11.190,0.030) -- (11.380,0.030) -- (11.380,0.468) -- (11.570,0.468) -- (11.570,1.216) -- (11.760,1.216) -- (11.760,1.609) -- (11.950,1.609) -- (11.950,1.632) -- (12.140,1.632) -- (12.140,1.385) -- (12.330,1.385) -- (12.330,1.056) -- (12.520,1.056) -- (12.520,0.854) -- (12.710,0.854) -- (12.710,0.598) -- (12.900,0.598) -- (12.900,0.478) -- (13.090,0.478) -- (13.090,0.343) -- (13.280,0.343) -- (13.280,0.248) -- (13.470,0.248) -- (13.470,0.173) -- (13.660,0.173) -- (13.660,0.150) -- (13.850,0.150) -- (13.850,0.096) -- (14.040,0.096) -- (14.040,0.079) -- (14.230,0.079) -- (14.230,0.054) -- (14.420,0.054) -- (14.420,0.035) -- (14.610,0.035) -- (14.610,0.028) -- (14.800,0.028) -- (14.800,0.021) -- (14.990,0.021) -- (14.990,0.012) -- (15.180,0.012) -- (15.180,0.015) -- (15.370,0.015) -- (15.370,0.007) -- (15.560,0.007) -- (15.560,0.007) -- (15.750,0.007) -- (15.750,0);
\draw[cExact,thick,fill=cExact,fill opacity=0.12] (11.000,0) -- (11.000,0.000) -- (11.190,0.000) -- (11.190,0.055) -- (11.380,0.055) -- (11.380,0.702) -- (11.570,0.702) -- (11.570,1.708) -- (11.760,1.708) -- (11.760,2.130) -- (11.950,2.130) -- (11.950,1.896) -- (12.140,1.896) -- (12.140,1.443) -- (12.330,1.443) -- (12.330,0.949) -- (12.520,0.949) -- (12.520,0.647) -- (12.710,0.647) -- (12.710,0.396) -- (12.900,0.396) -- (12.900,0.256) -- (13.090,0.256) -- (13.090,0.156) -- (13.280,0.156) -- (13.280,0.102) -- (13.470,0.102) -- (13.470,0.064) -- (13.660,0.064) -- (13.660,0.042) -- (13.850,0.042) -- (13.850,0.020) -- (14.040,0.020) -- (14.040,0.014) -- (14.230,0.014) -- (14.230,0.012) -- (14.420,0.012) -- (14.420,0.005) -- (14.610,0.005) -- (14.610,0.005) -- (14.800,0.005) -- (14.800,0.002) -- (14.990,0.002) -- (14.990,0.003) -- (15.180,0.003) -- (15.180,0.002) -- (15.370,0.002) -- (15.370,0.001) -- (15.560,0.001) -- (15.560,0.000) -- (15.750,0.000) -- (15.750,0);
\draw[->] (11.000,0) -- (15.900,0);
\draw (11.000,0) -- (11.000,2.300);
\draw (11.000,0) -- (11.000,-0.08) node[below,font=\scriptsize] {$0$};
\draw (12.188,0) -- (12.188,-0.08) node[below,font=\scriptsize] {$1$};
\draw (13.375,0) -- (13.375,-0.08) node[below,font=\scriptsize] {$2$};
\draw (14.562,0) -- (14.562,-0.08) node[below,font=\scriptsize] {$3$};
\draw (15.750,0) -- (15.750,-0.08) node[below,font=\scriptsize] {$4$};
\draw[gray,thick,dashed] (12.188,0) -- (12.188,2.300);
\draw[black,densely dotted,thick] (12.484,0) -- (12.484,2.300);
\node[anchor=north east,font=\scriptsize,align=right] at (15.750,2.300) {$d=3$\\20000 條};
\node[font=\scriptsize] at (7.875,-0.75) {$\Lambda_k\big/\frac{N+1}{\pi^2k^2}$};
\end{tikzpicture}}

\newcommand{\SixTabOneDee}{%
{\setlength{\tabcolsep}{3.5pt}\begin{tabular}{@{}lccccc@{}}
\toprule
$d$（條數） & 平均，$k=1,2,3$ & (6.61) & 標準差 & 平均$\div$(6.61) & 變異數$\div$(6.62)\\
\midrule
100（2000） & 0.1027, 0.0254, 0.0114 & 0.1021, 0.0255, 0.0113 & 0.0142, 0.0035, 0.0015 & 1.01, 1.00, 1.00 & 0.98, 0.97, 0.94\\
20（2000） & 0.1059, 0.0264, 0.0117 & 0.1051, 0.0263, 0.0117 & 0.0324, 0.0075, 0.0029 & 1.01, 1.01, 1.00 & 1.02, 0.87, 0.68\\
3（20000） & 0.1280, 0.0286, 0.0105 & 0.1267, 0.0317, 0.0141 & 0.0830, 0.0142, 0.0045 & 1.01, 0.90, 0.75 & 1.01, 0.47, 0.23\\
\bottomrule
\end{tabular}}}

\newcommand{\SixFigFixman}{%
\begin{tikzpicture}[>=Stealth]
\draw[cBar!70,fill=cBar!20] (0.000,0) rectangle (0.220,0.028);
\draw[cBar!70,fill=cBar!20] (0.220,0) rectangle (0.440,0.989);
\draw[cBar!70,fill=cBar!20] (0.440,0) rectangle (0.660,2.213);
\draw[cBar!70,fill=cBar!20] (0.660,0) rectangle (0.880,2.508);
\draw[cBar!70,fill=cBar!20] (0.880,0) rectangle (1.100,2.321);
\draw[cBar!70,fill=cBar!20] (1.100,0) rectangle (1.320,1.920);
\draw[cBar!70,fill=cBar!20] (1.320,0) rectangle (1.540,1.603);
\draw[cBar!70,fill=cBar!20] (1.540,0) rectangle (1.760,1.305);
\draw[cBar!70,fill=cBar!20] (1.760,0) rectangle (1.980,1.031);
\draw[cBar!70,fill=cBar!20] (1.980,0) rectangle (2.200,0.798);
\draw[cBar!70,fill=cBar!20] (2.200,0) rectangle (2.420,0.619);
\draw[cBar!70,fill=cBar!20] (2.420,0) rectangle (2.640,0.528);
\draw[cBar!70,fill=cBar!20] (2.640,0) rectangle (2.860,0.380);
\draw[cBar!70,fill=cBar!20] (2.860,0) rectangle (3.080,0.276);
\draw[cBar!70,fill=cBar!20] (3.080,0) rectangle (3.300,0.220);
\draw[cBar!70,fill=cBar!20] (3.300,0) rectangle (3.520,0.178);
\draw[cBar!70,fill=cBar!20] (3.520,0) rectangle (3.740,0.134);
\draw[cBar!70,fill=cBar!20] (3.740,0) rectangle (3.960,0.099);
\draw[cBar!70,fill=cBar!20] (3.960,0) rectangle (4.180,0.087);
\draw[cBar!70,fill=cBar!20] (4.180,0) rectangle (4.400,0.059);
\draw[cBar!70,fill=cBar!20] (4.400,0) rectangle (4.620,0.052);
\draw[cBar!70,fill=cBar!20] (4.620,0) rectangle (4.840,0.026);
\draw[cBar!70,fill=cBar!20] (4.840,0) rectangle (5.060,0.024);
\draw[cBar!70,fill=cBar!20] (5.060,0) rectangle (5.280,0.015);
\draw[cBar!70,fill=cBar!20] (5.280,0) rectangle (5.500,0.026);
\draw[cBar!70,fill=cBar!20] (5.500,0) rectangle (5.720,0.014);
\draw[cBar!70,fill=cBar!20] (5.720,0) rectangle (5.940,0.009);
\draw[cBar!70,fill=cBar!20] (5.940,0) rectangle (6.160,0.006);
\draw[cBar!70,fill=cBar!20] (6.160,0) rectangle (6.380,0.010);
\draw[cBar!70,fill=cBar!20] (6.380,0) rectangle (6.600,0.004);
\draw[cGauss,very thick] (0.000,0.000) -- (0.023,0.992) -- (0.046,1.361) -- (0.068,1.617) -- (0.091,1.811) -- (0.114,1.963) -- (0.137,2.086) -- (0.159,2.185) -- (0.182,2.265) -- (0.205,2.330) -- (0.228,2.382) -- (0.250,2.423) -- (0.273,2.454) -- (0.296,2.478) -- (0.319,2.494) -- (0.341,2.503) -- (0.364,2.507) -- (0.387,2.507) -- (0.410,2.501) -- (0.432,2.492) -- (0.455,2.480) -- (0.478,2.465) -- (0.501,2.447) -- (0.523,2.426) -- (0.546,2.403) -- (0.569,2.379) -- (0.592,2.353) -- (0.614,2.325) -- (0.637,2.297) -- (0.660,2.267) -- (0.682,2.237) -- (0.748,2.143) -- (0.815,2.045) -- (0.881,1.944) -- (0.948,1.844) -- (1.014,1.744) -- (1.081,1.646) -- (1.147,1.551) -- (1.214,1.459) -- (1.280,1.370) -- (1.347,1.285) -- (1.413,1.203) -- (1.480,1.126) -- (1.546,1.052) -- (1.613,0.983) -- (1.679,0.917) -- (1.746,0.855) -- (1.812,0.797) -- (1.879,0.742) -- (1.945,0.690) -- (2.012,0.642) -- (2.078,0.596) -- (2.145,0.554) -- (2.211,0.514) -- (2.278,0.477) -- (2.344,0.443) -- (2.411,0.411) -- (2.477,0.381) -- (2.544,0.353) -- (2.610,0.327) -- (2.677,0.302) -- (2.743,0.280) -- (2.810,0.259) -- (2.876,0.240) -- (2.943,0.222) -- (3.009,0.205) -- (3.076,0.189) -- (3.142,0.175) -- (3.209,0.162) -- (3.275,0.150) -- (3.342,0.138) -- (3.408,0.128) -- (3.475,0.118) -- (3.541,0.109) -- (3.608,0.100) -- (3.674,0.093) -- (3.741,0.085) -- (3.807,0.079) -- (3.874,0.073) -- (3.940,0.067) -- (4.007,0.062) -- (4.073,0.057) -- (4.140,0.053) -- (4.206,0.048) -- (4.273,0.045) -- (4.339,0.041) -- (4.406,0.038) -- (4.472,0.035) -- (4.539,0.032) -- (4.605,0.030) -- (4.672,0.027) -- (4.738,0.025) -- (4.805,0.023) -- (4.871,0.021) -- (4.938,0.020) -- (5.004,0.018) -- (5.071,0.017) -- (5.137,0.015) -- (5.204,0.014) -- (5.270,0.013) -- (5.337,0.012) -- (5.403,0.011) -- (5.470,0.010) -- (5.536,0.009) -- (5.603,0.009) -- (5.669,0.008) -- (5.736,0.007) -- (5.802,0.007) -- (5.869,0.006) -- (5.935,0.006) -- (6.002,0.005) -- (6.068,0.005) -- (6.135,0.004) -- (6.201,0.004) -- (6.268,0.004) -- (6.334,0.003) -- (6.401,0.003) -- (6.467,0.003) -- (6.534,0.003) -- (6.600,0.002);
\draw[cWalk,thick,dashed] (1.408,0) -- (1.408,2.800) node[above,font=\scriptsize,inner sep=1pt] {模擬平均 0.1280};
\draw[->] (0.000,0) -- (6.800,0);
\draw[->] (0.000,0) -- (0.000,3.050) node[above,font=\scriptsize] {密度};
\draw (0.000,0) -- (0.000,-0.07) node[below,font=\scriptsize] {$0$};
\draw (1.100,0) -- (1.100,-0.07) node[below,font=\scriptsize] {$0.1$};
\draw (2.200,0) -- (2.200,-0.07) node[below,font=\scriptsize] {$0.2$};
\draw (3.300,0) -- (3.300,-0.07) node[below,font=\scriptsize] {$0.3$};
\draw (4.400,0) -- (4.400,-0.07) node[below,font=\scriptsize] {$0.4$};
\draw (5.500,0) -- (5.500,-0.07) node[below,font=\scriptsize] {$0.5$};
\draw (6.600,0) -- (6.600,-0.07) node[below,font=\scriptsize] {$0.6$};
\draw (0.000,0.700) -- (-0.070,0.700) node[left,font=\scriptsize] {$2$};
\draw (0.000,1.400) -- (-0.070,1.400) node[left,font=\scriptsize] {$4$};
\draw (0.000,2.100) -- (-0.070,2.100) node[left,font=\scriptsize] {$6$};
\node[font=\scriptsize,anchor=west] at (2.4,1.9) {\color{cGauss}Gamma (6.73)：峰 0.0338、高 7.16};
\node[font=\scriptsize,anchor=west] at (2.4,1.45) {\color{cBar!80!black}模擬 $R_1^2/(N+1)$};
\draw[cBar!70,fill=cBar!20] (8.120,0) rectangle (8.340,0.032);
\draw[cBar!70,fill=cBar!20] (8.340,0) rectangle (8.560,0.495);
\draw[cBar!70,fill=cBar!20] (8.560,0) rectangle (8.780,1.351);
\draw[cBar!70,fill=cBar!20] (8.780,0) rectangle (9.000,1.958);
\draw[cBar!70,fill=cBar!20] (9.000,0) rectangle (9.220,2.138);
\draw[cBar!70,fill=cBar!20] (9.220,0) rectangle (9.440,2.018);
\draw[cBar!70,fill=cBar!20] (9.440,0) rectangle (9.660,1.817);
\draw[cBar!70,fill=cBar!20] (9.660,0) rectangle (9.880,1.587);
\draw[cBar!70,fill=cBar!20] (9.880,0) rectangle (10.100,1.272);
\draw[cBar!70,fill=cBar!20] (10.100,0) rectangle (10.320,1.063);
\draw[cBar!70,fill=cBar!20] (10.320,0) rectangle (10.540,0.837);
\draw[cBar!70,fill=cBar!20] (10.540,0) rectangle (10.760,0.646);
\draw[cBar!70,fill=cBar!20] (10.760,0) rectangle (10.980,0.528);
\draw[cBar!70,fill=cBar!20] (10.980,0) rectangle (11.200,0.423);
\draw[cBar!70,fill=cBar!20] (11.200,0) rectangle (11.420,0.305);
\draw[cBar!70,fill=cBar!20] (11.420,0) rectangle (11.640,0.242);
\draw[cBar!70,fill=cBar!20] (11.640,0) rectangle (11.860,0.165);
\draw[cBar!70,fill=cBar!20] (11.860,0) rectangle (12.080,0.144);
\draw[cBar!70,fill=cBar!20] (12.080,0) rectangle (12.300,0.108);
\draw[cBar!70,fill=cBar!20] (12.300,0) rectangle (12.520,0.081);
\draw[cBar!70,fill=cBar!20] (12.520,0) rectangle (12.740,0.070);
\draw[cBar!70,fill=cBar!20] (12.740,0) rectangle (12.960,0.053);
\draw[cBar!70,fill=cBar!20] (12.960,0) rectangle (13.180,0.027);
\draw[cBar!70,fill=cBar!20] (13.180,0) rectangle (13.400,0.025);
\draw[cBar!70,fill=cBar!20] (13.400,0) rectangle (13.620,0.028);
\draw[cBar!70,fill=cBar!20] (13.620,0) rectangle (13.840,0.017);
\draw[cBar!70,fill=cBar!20] (13.840,0) rectangle (14.060,0.018);
\draw[cBar!70,fill=cBar!20] (14.060,0) rectangle (14.280,0.015);
\draw[cBar!70,fill=cBar!20] (14.280,0) rectangle (14.500,0.006);
\draw[cExact,very thick] (7.933,-0.000) -- (7.955,-0.000) -- (7.976,-0.000) -- (7.998,-0.000) -- (8.019,-0.000) -- (8.041,-0.000) -- (8.063,0.000) -- (8.084,0.000) -- (8.106,0.000) -- (8.128,0.000) -- (8.149,0.001) -- (8.171,0.002) -- (8.192,0.004) -- (8.214,0.010) -- (8.236,0.018) -- (8.257,0.031) -- (8.279,0.051) -- (8.301,0.077) -- (8.322,0.110) -- (8.344,0.152) -- (8.365,0.202) -- (8.387,0.260) -- (8.409,0.325) -- (8.430,0.397) -- (8.452,0.475) -- (8.474,0.558) -- (8.495,0.645) -- (8.517,0.734) -- (8.538,0.826) -- (8.560,0.918) -- (8.582,1.012) -- (8.648,1.286) -- (8.715,1.533) -- (8.781,1.742) -- (8.848,1.907) -- (8.914,2.029) -- (8.981,2.111) -- (9.047,2.157) -- (9.114,2.171) -- (9.180,2.160) -- (9.247,2.129) -- (9.313,2.080) -- (9.380,2.019) -- (9.446,1.949) -- (9.513,1.871) -- (9.579,1.790) -- (9.646,1.705) -- (9.712,1.620) -- (9.779,1.535) -- (9.845,1.450) -- (9.912,1.368) -- (9.978,1.288) -- (10.045,1.210) -- (10.111,1.136) -- (10.178,1.064) -- (10.244,0.996) -- (10.311,0.932) -- (10.377,0.870) -- (10.444,0.812) -- (10.510,0.757) -- (10.577,0.706) -- (10.643,0.657) -- (10.710,0.611) -- (10.776,0.569) -- (10.843,0.529) -- (10.909,0.491) -- (10.976,0.456) -- (11.042,0.423) -- (11.109,0.392) -- (11.175,0.364) -- (11.242,0.337) -- (11.308,0.313) -- (11.375,0.290) -- (11.441,0.268) -- (11.508,0.248) -- (11.574,0.230) -- (11.641,0.212) -- (11.707,0.196) -- (11.774,0.182) -- (11.840,0.168) -- (11.907,0.155) -- (11.973,0.143) -- (12.040,0.132) -- (12.106,0.122) -- (12.173,0.113) -- (12.239,0.104) -- (12.306,0.096) -- (12.372,0.089) -- (12.439,0.082) -- (12.505,0.076) -- (12.572,0.070) -- (12.638,0.064) -- (12.705,0.059) -- (12.771,0.055) -- (12.838,0.050) -- (12.904,0.047) -- (12.971,0.043) -- (13.037,0.040) -- (13.104,0.036) -- (13.170,0.034) -- (13.237,0.031) -- (13.303,0.028) -- (13.370,0.026) -- (13.436,0.024) -- (13.503,0.022) -- (13.569,0.020) -- (13.636,0.019) -- (13.702,0.017) -- (13.769,0.016) -- (13.835,0.015) -- (13.902,0.014) -- (13.968,0.012) -- (14.035,0.011) -- (14.101,0.011) -- (14.168,0.010) -- (14.234,0.009) -- (14.301,0.008) -- (14.367,0.008) -- (14.434,0.007) -- (14.500,0.006);
\draw[cWalk,thick,dashed] (9.733,0) -- (9.733,2.800) node[above,font=\scriptsize,inner sep=1pt] {平均 $\frac16$};
\draw[->] (7.900,0) -- (14.700,0);
\draw[->] (7.900,0) -- (7.900,3.050) node[above,font=\scriptsize] {密度};
\draw (7.900,0) -- (7.900,-0.07) node[below,font=\scriptsize] {$0$};
\draw (9.000,0) -- (9.000,-0.07) node[below,font=\scriptsize] {$0.1$};
\draw (10.100,0) -- (10.100,-0.07) node[below,font=\scriptsize] {$0.2$};
\draw (11.200,0) -- (11.200,-0.07) node[below,font=\scriptsize] {$0.3$};
\draw (12.300,0) -- (12.300,-0.07) node[below,font=\scriptsize] {$0.4$};
\draw (13.400,0) -- (13.400,-0.07) node[below,font=\scriptsize] {$0.5$};
\draw (14.500,0) -- (14.500,-0.07) node[below,font=\scriptsize] {$0.6$};
\draw (7.900,0.700) -- (7.830,0.700) node[left,font=\scriptsize] {$2$};
\draw (7.900,1.400) -- (7.830,1.400) node[left,font=\scriptsize] {$4$};
\draw (7.900,2.100) -- (7.830,2.100) node[left,font=\scriptsize] {$6$};
\node[font=\scriptsize,anchor=west] at (10.500,1.9) {\color{cExact}Fixman (6.74)：峰 0.111、高 6.20};
\node[font=\scriptsize,anchor=west] at (10.500,1.45) {\color{cBar!80!black}模擬 $\mathrm{Tr}\,\mathbf T/(N+1)$};
\end{tikzpicture}}

\newcommand{\SixTabErrata}{%
\begin{tabular}{@{}lll@{}}
\toprule
位置 & 書上 & 應為\\
\midrule
p.~135，(6.16) & {\color{cWrong}$(d-1)\sum_{k=1}^{N}R_k^2$} & {\color{cExact}$(d-1)\big(\sum_{k=1}^{d}R_k^2\big)^2$}\\
p.~138，練習 6.3(b)(c) & {\color{cWrong}(b) $j<n$／$j\ge n$ 配 (c) $\sum_{j=1}^{N+1}$} & {\color{cExact}點編號 $j=0..N$：$\sum_{j=0}^{N}$}\\
p.~139，(6.22) & {\color{cWrong}$\alpha<\beta$；$\alpha>\beta$} & {\color{cExact}$\alpha\le\beta$；$\alpha>\beta$（含對角線）}\\
p.~139，(6.23) & {\color{cWrong}$a_{x_1,x_2}=x_<(N+1-x_>)$} & {\color{cExact}$a_{x_1,x_2}=\frac{x_<(N+1-x_>)}{(N+1)^2}$}\\
p.~141，(6.33a) & {\color{cWrong}$\sum_{\alpha,\gamma}a_{\alpha,\alpha}a_{\delta,\delta}$} & {\color{cExact}$\sum_{\alpha,\delta}a_{\alpha,\alpha}a_{\delta,\delta}$}\\
p.~141，(6.35a) & {\color{cWrong}$\frac{1}{d^2}\frac{(N+1)^2}{\pi^4d^2}\big[\big(\sum_n\frac1{n^2}\big)+\cdots\big]$} & {\color{cExact}$\frac{1}{d^2}\frac{(N+1)^2}{\pi^4}\big[\big(\sum_n\frac1{n^2}\big)^2+\cdots\big]$}\\
p.~141，(6.35b) & {\color{cWrong}$\frac{\pi^4}{36}+\frac{\pi^4}{90}$} & {\color{cExact}$\frac{\pi^4}{36}+\frac{2\pi^4}{90}$}\\
p.~146，(6.52) & {\color{cWrong}$\sum_i\delta(\lambda-\lambda_i)$} & {\color{cExact}$\pi\sum_i\delta(\lambda-\lambda_i)$}\\
p.~146，(6.54) & {\color{cWrong}$\langle R_1^2\rangle:\langle R_2^2\rangle:\langle R_3^2\rangle=13.76:3.03:1$} & {\color{cExact}$\approx11.9:2.7:1$（書上的數接近比值的平均）}\\
p.~153，(6.76) & {\color{cWrong}$K(R)=\prod\big(1-\frac{2is\alpha_n}{d}\big)^{-d/2}$} & {\color{cExact}$K(s)=\prod\big(1+\frac{2is\alpha_n}{d}\big)^{-d/2}$}\\
p.~159，(6:S2-1) & {\color{cWrong}$\cdots\big/\big(\frac{\langle T_{11}^2\rangle}{\langle T_{11}T_{22}\rangle}+d-1\big)$} & {\color{cExact}分母再乘 $\frac{\langle T_{11}T_{22}\rangle}{\langle T_{12}^2\rangle}$}\\
p.~164，(6:S2-25) 前 & {\color{cWrong}Adding (6:S2-16), (6:S2-21), and (6:S2-23)} & {\color{cExact}(6:S2-24)}\\
\bottomrule
\end{tabular}}


\section{漫步的形狀（第 6 章）}

\begin{keybox}[先看這裡]
前五章只問醉漢「停在哪裡」。這一章問：他走過的那一串足跡\textbf{長什麼樣子}？
\textbf{平均起來是圓的}（1000 條疊起來是一團圓雲），\textbf{單獨一條卻幾乎總是長條}。描述一條漫步的形狀要用\textbf{迴轉半徑張量} $\mathbf T$（像轉動慣量的矩陣）：
特徵向量是延伸最長、最短的方向，特徵值 $R_1^2\ge R_2^2\ge\cdots$ 是那些方向的平方寬度。把它們合成一個數，就是\textbf{非球度}：0 是圓球、1 是一條線。
\textbf{結論}：$d$ 維的平均非球度 $A_d=\frac{4+2d}{4+5d}$，三維是 $\frac{10}{19}=0.5263$；維度比步數多得多時，每條漫步都是同一座「樓梯」，$R_k^2=\frac{N+1}{\pi^2k^2}$、非球度 $\frac25$；
有限的 $d$ 用 $1/d$ 展開：尖峰相對地往右移 $\frac{3}{4d}$、相對寬度 $\sqrt{2/d}$。總整理在 6.4 節。

\textbf{記號}：$N$ 步的漫步走過 $N+1$ 個點 $\vec r_0=0,\vec r_1,\dots,\vec r_N$；步伐 $\vec\eta_\alpha=\vec r_\alpha-\vec r_{\alpha-1}$，$\langle|\vec\eta|^2\rangle=l^2=1$；維度 $d$。
格子漫步＝每步沿 $2d$ 個方向之一走 1；高斯步伐＝每個分量是常態、變異數 $\frac1d$。Kramers 矩陣 $a_{\alpha\beta}$；$A_d$＝平均的比，$A'_d$＝個別值的平均。

\textbf{全章的主例}：4 步「東、東、北、西」：$\mathbf T=\frac1{25}\big(\begin{smallmatrix}14&3\\3&6\end{smallmatrix}\big)$、$R_1^2=\frac35$、$R_2^2=\frac15$、非球度 $\frac14$；
固定種子的二維 1000 步（書上 Fig.~6.4 的樣子）：$R_1^2=258.8$、$R_2^2=31.6$；$N=3$ 的樓梯：$\mathbf T=\frac1{16}\big(\begin{smallmatrix}3&2&1\\2&4&2\\1&2&3\end{smallmatrix}\big)$。
\begin{center}
\small
\begin{tabular}{@{}lll@{\qquad}lll@{}}
\toprule
節 & 內容 & 動畫 & 節 & 內容 & 動畫\\
\midrule
6.1--6.1.1 & 一團雲、一條長條 & W601 & 6.2 & $d\gg N$：樓梯 & W604\\
6.1.2--6.1.4 & 迴轉半徑張量、主軸 & W602 & 6.2.1 & Kramers 恆等式、配對 & W605\\
6.1.4 & 非球度 & W603 & 6.2.2--6.2.5 & $1/d$ 展開 & W606\\
 & & & 6.2.6--6.3 & 分布、Fixman、自迴避 & W607\\
\bottomrule
\end{tabular}
\end{center}
（上表是書上的節號。）數字一律由 \texttt{codes/utils/walk\_shape.py} 重算，圖表由 \texttt{ch6/make\_figs.py} 產生並先做數字檢查；逐式校對見 \texttt{ref/md/ch06\_formulas.md}。
連續播放：\texttt{Ch6Demo}（W601--W607）。
\end{keybox}

% =====================================================================
\subsection{迴轉半徑張量與非球度}
\leavevmode\bk{pp.~127--135}\quad 終點只是一個點；這一節要描述整串足跡的形狀，而且要用\textbf{轉一轉也不變}的量。

\subsubsection{一團雲與一條長條}
\leavevmode\bk{pp.~127--129}\ \anim{W601}\par\nopagebreak

\textbf{問}：1000 個醉漢各走 100 步（二維），疊起來是一團圓雲（書上 Fig.~6.1）。單獨一條也是圓的嗎？先猜。

\textbf{小例子}：4 步「東、東、北、西」走過 5 個點 $(0,0),(1,0),(2,0),(2,1),(1,1)$。重心是這些點的平均 $(\frac65,\frac25)$；
五個點到重心的平方距離是 $\frac85,\frac15,\frac45,1,\frac25$，平均起來
\[R_g^2=\frac15\Big(\frac85+\frac15+\frac45+1+\frac25\Big)=\frac15\cdot\frac{20}{5}=\frac45 .\]
這就是\textbf{迴轉半徑的平方}：每個點當成質量 1，$R_g^2$ 是繞重心的轉動慣量除以總質量（圖 \ref{six-fig-shape} 左）。

\textbf{一大群}：種子 601 的 1000 條 100 步。把每一條的張量 $\mathbf T$（下一小節）平均起來，得 $\bar T_{11}=8.35$、$\bar T_{22}=8.59$、$\bar T_{12}=0.03$
（理論值 $\frac{N(N+2)}{12(N+1)}=8.42$、$8.42$、$0$），兩個特徵值 8.60、8.34：平均起來幾乎是圓的。
\textbf{揭曉}：單獨一條卻幾乎都是長條。書上 Fig.~6.3 那種 6 條 1000 步（種子 603）的 $R_1^2/R_2^2$ 是 2.1、3.0、6.7、5.5、14.5、4.6；
上面那 1000 條 100 步的平均 $\langle R_1^2\rangle=14.14$、$\langle R_2^2\rangle=2.80$，比 5.05。雲是圓的，只是因為每一條往不同的方向長；
要描述「一條」的形狀，得用不隨旋轉改變的量。

\subsubsection{迴轉半徑張量：形狀的指紋}
\leavevmode\bk{pp.~129--133}\ \anim{W602}\par\nopagebreak

\textbf{問}：怎麼用幾個數字說出「這條漫步往哪個方向長、長多少」？提示：轉動慣量。

\textbf{定義}（$N+1$ 個點、除以 $N+1$）：
\[T_{kl}=\frac{1}{N+1}\sum_{j=0}^{N}\big(r_{jk}-\langle r_k\rangle\big)\big(r_{jl}-\langle r_l\rangle\big),\qquad
\langle r_k\rangle=\frac{1}{N+1}\sum_{j=0}^{N}r_{jk}.\tag{6.1, 6.2}\]
二維就是 $2\times2$ 的共變異矩陣 (6.3)：對角線是 $x$、$y$ 方向的平方寬度，非對角線量 $x$ 與 $y$ 一起變的程度；$\mathrm{Tr}\,\mathbf T=R_g^2$。
〔注〕書上 (6.1)(6.2) 寫「$N$ 個足跡、除以 $N$」，(6.22) 起才改成「$N$ 步、$N+1$ 個點」；本講義一律用後者（$N$ 大時沒有差別）。

\textbf{算到底}：小例子對重心的座標是 $x-\frac65=-\frac65,-\frac15,\frac45,\frac45,-\frac15$、$y-\frac25=-\frac25,-\frac25,-\frac25,\frac35,\frac35$；
乘積的分母都是 25，再除以 5 個點：
\[T_{11}=\tfrac{36+1+16+16+1}{125}=\frac{14}{25},\qquad T_{22}=\tfrac{4+4+4+9+9}{125}=\frac{6}{25},\qquad
T_{12}=\tfrac{12+2-8+12-3}{125}=\frac{3}{25}.\]
$\big(\begin{smallmatrix}14&3\\3&6\end{smallmatrix}\big)$ 的特徵值是 $\frac{20\pm\sqrt{8^2+4\cdot3^2}}{2}=15,\ 5$，對角化 $\mathbf T=\mathrm{diag}(R_1^2,R_2^2)$：
\[R_1^2=\frac{15}{25}=\frac35,\qquad R_2^2=\frac15,\qquad\text{主軸 }(3,1)\ \ \big(\arctan\tfrac13=18.43^\circ\big).\tag{6.4}\]
不變量 (6.6)--(6.9)：$\mathrm{Tr}\,\mathbf T=\frac45$、$\mathrm{Tr}\,\mathbf T^2=\frac{9+1}{25}=\frac25$、$\det\mathbf T=\frac3{25}=\frac35\cdot\frac15$
（二維 $\det\mathbf T=\frac12[(\mathrm{Tr}\,\mathbf T)^2-\mathrm{Tr}\,\mathbf T^2]$；三維 $\det\mathbf T=\frac16[(\mathrm{Tr}\,\mathbf T)^3+2\,\mathrm{Tr}\,\mathbf T^3-3\,\mathrm{Tr}\,\mathbf T\,\mathrm{Tr}\,\mathbf T^2]$）。
它們都只用特徵值寫得出來，所以轉一轉不會變。

\textbf{特徵向量是什麼方向？}（補充 1）沿單位向量 $\vec n=(\cos\varphi,\sin\varphi)$ 的平方寬度是 $\vec n\cdot\mathbf T\cdot\vec n$。小例子在 $\varphi=18.43^\circ$ 時最大，是 $\frac35$；
$\varphi=108.43^\circ$ 時最小，是 $\frac15$。所以特徵向量＝延伸最長、最短的方向，特徵值＝那個方向的平方寬度。
\textbf{旋轉}：$\mathbf T'=\mathbf R\,\mathbf T\,\mathbf R^{-1}$（(6:S1-17)），矩陣會變、特徵值不變。對所有方向平均，$\langle\mathbf T\rangle=\bar T\,\mathbf I$（(6.5)）——形狀的資訊被抹掉，這就是那團圓雲。

\textbf{主例}（書上 Fig.~6.4 的樣子；種子 643 的二維 1000 步，圖 \ref{six-fig-shape} 右）：重心 $(-17.9,-8.9)$，
$T_{11}=207.5$、$T_{22}=82.9$、$T_{12}=95.0$；$R_1^2=258.8$、$R_2^2=31.6$（$R_1=16.09$、$R_2=5.62$，$R_1^2/R_2^2=8.18$），主軸 $28.4^\circ$。
整條繞重心轉 $30^\circ$：$T_{11}=94.1$、$T_{22}=196.3$、$T_{12}=101.4$，特徵值仍是 258.8、31.6，主軸變成 $58.4^\circ$。
〔注〕Fig.~6.4 說兩條線的長度「正比於特徵值」；本講義與動畫畫成半長 $2R_i$（長度的單位）。

\begin{figure}[htbp]
\centering
\SixFigSmall\hfill\SixFigMain\\[6pt]
{\small\SixTabMain}
\caption{迴轉半徑張量。左上：「東、東、北、西」的 5 個點（數字是到重心的平方距離）、重心（十字）、兩條主軸（半長 $2R_i$）與半軸 $2R_1$、$2R_2$ 的橢圓。
右上：種子 643 的二維 1000 步（書上 Fig.~6.4 的樣子；$\circ$ 是起點，比例尺 10 步）。下：數字；轉 $30^\circ$ 之後矩陣全變了，特徵值與非球度不變}\label{six-fig-shape}
\end{figure}

\subsubsection{非球度：0 是圓球、1 是一條線}
\leavevmode\bk{pp.~133--135}\ \anim{W603}\par\nopagebreak

\textbf{問}：三維漫步平均起來有多扁？在 0（圓球）到 1（一條線）之間猜一個數。

\textbf{一條的非球度}：特徵值兩兩相減的平方和，用不變量寫就不必對角化（三維：$3\,\mathrm{Tr}\,\mathbf T^2-(\mathrm{Tr}\,\mathbf T)^2$）：
\[\sum_{i<j}(R_i^2-R_j^2)^2=d\,\mathrm{Tr}\,\mathbf T^2-(\mathrm{Tr}\,\mathbf T)^2,\qquad
\hat A=\frac{\sum_{i<j}(R_i^2-R_j^2)^2}{(d-1)\big(\sum_iR_i^2\big)^2}\tag{6.10, 6.13}\]
（$\hat A$ 是本講義的記號）。分母的 $d-1$ 讓「一條線」剛好是 1。\textbf{例子}：圓球 $(1,1,1)$：0；線 $(1,0,0)$：$\frac{1+1+0}{2\cdot1}=1$；
雪茄 $(4,1,1)$：$\frac{9+9+0}{2\cdot36}=\frac14$；扁盤 $(1,1,0)$：$\frac{0+1+1}{2\cdot4}=\frac14$——雪茄和扁盤分不出來，一個數只能說「有多不圓」。
二維小例子 $(\frac35,\frac15)$：$\frac{(2/5)^2}{(4/5)^2}=\frac14$；主例 0.6117。

\textbf{平均非球度}：書上把分子、分母\textbf{各自平均再相除}（平均的比），再用 $\mathbf T$ 的分量寫（補充 2）：
\[A_d=\frac{\sum_{i<j}\langle(R_i^2-R_j^2)^2\rangle}{(d-1)\big\langle(\sum_iR_i^2)^2\big\rangle}
=\frac{\langle T_{11}^2\rangle-\langle T_{11}T_{22}\rangle+d\langle T_{12}^2\rangle}{\langle T_{11}^2\rangle+(d-1)\langle T_{11}T_{22}\rangle}
=\frac{4+2d}{4+5d}.\tag{6.12, 6.14, 6.15}\]
最後一步要知道三個平均，\ref{six-sec-kramers} 節用高斯步伐算出來。\textbf{揭曉}：三維 $A_3=\frac{10}{19}=0.5263$——大約「一半扁」（圖 \ref{six-fig-a} 上）。

\textbf{模擬}（書上 Fig.~6.5 的樣子；三維格子 2 萬條 100 步，種子 605，圖 \ref{six-fig-a} 下）：平均的比 0.5231（$N=100$ 格子漫步的精確值 0.5235，$N\to\infty$ 是 0.5263）；
但個別的 $\hat A$ 從 0 到 1 都有，小於 0.5 的占 69.3\%：平均不代表典型。
\textbf{個別值的平均}（(6.16)--(6.18)）：
\[A'_d=\Big\langle\frac{\sum_{i<j}(R_i^2-R_j^2)^2}{(d-1)\big(\sum_{k=1}^{d}R_k^2\big)^2}\Big\rangle=\frac d4\Big[3+\frac4d-\frac d2M_{d/2}\Big],\qquad M_p=\int_0^\infty x^{p+1}\sinh^{-p}x\,dx .\]
$A'_2=\frac12(5-\frac72\zeta(3))=0.396400$、$A'_3=0.394274$、$A'_4=4-3\zeta(3)=0.393829$；模擬的 $\hat A$ 平均是 0.3939。
比 $\frac{10}{19}$ 小很多：$A_d$ 其實是用 $(\mathrm{Tr}\,\mathbf T)^2$ 加權的平均，張得開的漫步權重大——兩者差這麼多，表示張得越開的漫步通常也越長條。
\begin{errbox}[p.~135，(6.16)]
書上分母寫成 $(d-1)\sum_{k=1}^{N}R_k^2$：沒有平方、上限是 $N$。應為 $(d-1)\big(\sum_{k=1}^{d}R_k^2\big)^2$。照書上的寫法，分子是長度的 4 次方、分母是 2 次方，
整條漫步放大 2 倍，小例子的值就從 0.2 變成 0.8——量的就不是形狀了。
\end{errbox}

\begin{figure}[htbp]
\centering
{\small\SixTabA}\\[6pt]
\SixFigHist
\caption{非球度。上：平均的比 $A_d$（(6.15)）與個別值的平均 $A'_d$（(6.17)）。下：三維格子 2 萬條 100 步的個別非球度 $\hat A$
（書上 Fig.~6.5 的樣子，20 格、寬 0.05）；虛線是平均的比，點線是個別值的平均}\label{six-fig-a}
\end{figure}

% =====================================================================
\subsection{很高的維度：樓梯、Kramers 恆等式與 \texorpdfstring{$1/d$}{1/d} 展開}
\leavevmode\bk{pp.~135--154}\quad $A_d$ 的最後一步要算 $\mathbf T$ 分量的平均。書上的辦法：先看維度很高的極限（每條漫步形狀都一樣），
再把 $\mathbf T$ 寫成步伐的二次式，用 $1/d$ 一階一階地修正。

\subsubsection{\texorpdfstring{$d\gg N$}{d >> N}：每條漫步都是同一座樓梯}
\leavevmode\bk{pp.~135--138}\ \anim{W604}\par\nopagebreak

\textbf{問}：空間有 100 維，醉漢走 10 步，會不會走回已經走過的方向（同一條軸）？

\textbf{數一數}（練習 6.2）：第一步有 $2d$ 種走法；要和前面都垂直，第二步只剩 $2d-2$ 種……所以每一步都換新軸的比例是
\[\frac{(2d)(2d-2)\cdots(2d-2N+2)}{(2d)^N}=\prod_{k<N}\Big(1-\frac kd\Big)\approx e^{-N^2/2d}\qquad(N\ll\sqrt d).\]
$d=100$、$N=10$：精確 0.6282，近似 $e^{-1/2}=0.6065$。$d\gg N^2$ 時幾乎每一步都垂直於之前所有步；轉一轉、照照鏡子，就能讓第 $n$ 步沿第 $n$ 軸走 $l$——
\textbf{每條漫步都是同一座樓梯}，形狀不再隨機。

\textbf{$N=3$ 的樓梯（算到底）}：點 $(0,0,0),(1,0,0),(1,1,0),(1,1,1)$。第 $n$ 個座標在前 $n$ 個點是 0、其餘是 1，所以重心 $\langle x_n\rangle=\frac{N+1-n}{N+1}=\frac34,\frac12,\frac14$（練習 6.3(a)）。
例如 $x_1=0,1,1,1$ 對平均的偏差是 $-\frac34,\frac14,\frac14,\frac14$，$T_{11}=\frac14(\frac9{16}+\frac3{16})=\frac3{16}$；同樣算出
\[\mathbf T=\frac1{16}\begin{pmatrix}3&2&1\\2&4&2\\1&2&3\end{pmatrix},\qquad R_k^2=\frac{2+\sqrt2}8,\ \frac18,\ \frac{2-\sqrt2}8=0.4268,\ 0.1250,\ 0.0732,\]
特徵向量 $(1,\sqrt2,1)$、$(1,0,-1)$、$(1,-\sqrt2,1)$：正好是 $\sin\frac{k\pi n}{4}$ 在 $n=1,2,3$ 的值。
一般的 $N$（練習 6.3）：$T_{nm}=\frac{\min(n,m)\,(N+1-\max(n,m))}{(N+1)^2}$，特徵向量精確是 $\sin\frac{k\pi n}{N+1}$；$N$ 大時
\[R_k^2=\frac{N+1}{\pi^2k^2},\qquad A_\infty=\frac{\sum_k(R_k^2)^2}{(\sum_kR_k^2)^2}=\frac{\sum k^{-4}}{(\sum k^{-2})^2}=\frac{\pi^4/90}{(\pi^2/6)^2}=\frac{1.0823}{2.7058}=\frac25\tag{6.19, 6.20}\]
（$d\to\infty$ 時 (6.13) 只剩 $\mathrm{Tr}\,\mathbf T^2$ 與 $(\mathrm{Tr}\,\mathbf T)^2$；$\sum k^{-2}=1.6449$）。$N=3$ 用大 $N$ 公式是 0.4053、0.1013、0.0450：最大的特徵值只差 5\%，小的差得多（表 \ref{six-tab-stair}）。
〔注〕(6.19) 是大 $N$ 的極限；精確值 $R_k^2=\frac{1}{4(N+1)\sin^2\frac{k\pi}{2(N+1)}}$。有限 $N$ 的樓梯非球度 $\frac{\mathrm{Tr}\,\mathbf T^2}{(\mathrm{Tr}\,\mathbf T)^2}$：
$N=3$ 是 $\frac{13}{25}=0.5200$，$N=10$ 是 0.4150，$N=100$ 是 0.4002 $\to\frac25$。
\begin{errbox}[p.~138，練習 6.3(b)(c)]
點的編號前後不一致：(b) 的 $j<n$／$j\ge n$ 是點 $j=0..N$ 的編號，(c) 的和要寫 $\sum_{j=0}^{N}$。照書上的 $\sum_{j=1}^{N+1}$ 配 (b)，平均會變成 $\frac{N+2-n}{N+1}$。
\end{errbox}

\textbf{和三維比}：(6.19) 的比是 $1:\frac14:\frac19$，即 $\langle R_1^2\rangle:\langle R_2^2\rangle:\langle R_3^2\rangle=9:2.25:1$（(6.53)）；三維格子 2 萬條 100 步的模擬是 11.95 : 2.69 : 1——
差不多，三維的漫步更長條一些。
\begin{errbox}[p.~146，(6.54)]
書上引 Bishop \& Michels 的模擬寫 $\langle R_1^2\rangle:\langle R_2^2\rangle:\langle R_3^2\rangle=13.76:3.03:1$，應為 $\approx11.9:2.7:1$（上面的模擬；
書上自己的 (6.69)--(6.71) 換算出 $1.01\cdot9/0.75=12.12$、$0.91\cdot2.25/0.75=2.73$，也一致）。
13.76 : 3.03 比較接近\textbf{比值的平均} $\langle R_1^2/R_3^2\rangle:\langle R_2^2/R_3^2\rangle$（模擬 13.69 : 2.98）：又是「平均的比 $\ne$ 比值的平均」。
\end{errbox}

\begin{table}[htbp]
\centering
\small
\SixTabStair
\caption{上：樓梯（$N=3$）與 Kramers 矩陣（$N=4$）的特徵值，精確值與大 $N$ 公式 $\frac{N+1}{\pi^2k^2}$（$k$ 小時接近）。
下左：$d\to\infty$ 的樓梯非球度 $\to\frac25$；下右：三維的主迴轉半徑比（模擬是三維格子 2 萬條 100 步）}\label{six-tab-stair}
\end{table}

\subsubsection{Kramers 恆等式：形狀是步伐的二次式}\label{six-sec-kramers}
\leavevmode\bk{pp.~138--142}\ \anim{W605}\par\nopagebreak

\textbf{問}：$\mathbf T$ 是 $N+1$ 個點的平均。能不能直接寫成「步伐乘步伐」？

\textbf{兩兩差}（練習 6.1）：$T_{kl}=\frac{1}{2(N+1)^2}\sum_i\sum_j(r_{ik}-r_{jk})(r_{il}-r_{jl})$，不必先算重心。點 $p<q$ 的差 $\vec r_q-\vec r_p=\sum_{\alpha=p+1}^{q}\vec\eta_\alpha$ 是步伐的和，展開就得
\[T_{ij}=\sum_{\alpha,\beta=1}^{N}a_{\alpha\beta}\,\eta_{\alpha,i}\eta_{\beta,j},\qquad
a_{\alpha\beta}=\frac{\min(\alpha,\beta)\,\big(N+1-\max(\alpha,\beta)\big)}{(N+1)^2}.\tag{6.21, 6.22}\]
\textbf{為什麼是這個數}（補充 3）：$\eta_m\eta_n$（$m\le n$）出現在點對 $p<q$ 的差裡，條件是 $p<m$、$q\ge n$：$p$ 有 $m$ 種、$q$ 有 $N+1-n$ 種，共 $m(N+1-n)$ 對。
\textbf{算到底}（$N=4$）：
\[\mathbf a=\frac1{25}\begin{pmatrix}4&3&2&1\\3&6&4&2\\2&4&6&3\\1&2&3&4\end{pmatrix},\qquad \mathrm{Tr}\,\mathbf a=\frac{20}{25}=\frac45 .\]
「東東北西」的 $x$ 分量 $(1,1,0,-1)$、$y$ 分量 $(0,0,1,0)$。$25\,\mathbf a$ 乘 $x$ 分量得 $(6,7,3,-1)$，所以 $T_{11}=\frac{6+7+0+1}{25}=\frac{14}{25}$、$T_{12}=\frac3{25}$，
$T_{22}=a_{33}=\frac6{25}$——和定義式一樣。樓梯的步伐是單位向量 $\vec e_\alpha$，代入立刻得 $T_{nm}=a_{nm}$。
\begin{errbox}[p.~139，(6.22)(6.23)]
(6.22) 只寫 $\alpha<\beta$、$\alpha>\beta$ 兩種情形，應為 $\alpha\le\beta$、$\alpha>\beta$：對角線 $a_{\alpha\alpha}=\frac{\alpha(N+1-\alpha)}{(N+1)^2}$ 不可少——少了它，「東東北西」會算出 $T_{11}=T_{22}=0$。
(6.23) 少了 $\frac1{(N+1)^2}$：$a_{x_1,x_2}=\frac{x_<(N+1-x_>)}{(N+1)^2}$，才接得上 (6.24) 的 $x(1-y)$（$x<y$）。
\end{errbox}

\textbf{正弦模態}（(6.24)--(6.27)）：令 $x=\frac{\alpha}{N+1}$，$a_{\alpha\beta}=x_<(1-x_>)$ 是 $-\frac{d^2}{dx^2}$（兩端為 0）的反算子，特徵函數 $\sin k\pi x$、特徵值 $\frac1{\pi^2k^2}$（練習 6.4）；
對 $\beta$ 求和約等於 $(N+1)\int dy$，所以 $\mathbf a$ 的特徵值 $\approx\frac{N+1}{\pi^2k^2}$，而特徵向量\textbf{精確}就是 $\sin\frac{k\pi\alpha}{N+1}$。
$N=4$：特徵值 0.5236、0.1447、0.0764、0.0553，大 $N$ 公式 0.5066、0.1267、0.0563、0.0317（表 \ref{six-tab-stair}）。〔注〕(6.26a) 的和寫到 $N+1$；那一項是 0，不影響。

\textbf{高斯步伐與配對}：每個分量獨立、常態、變異數 $\frac1d$（(6.28)--(6.30)）：
\[P(\eta_{\alpha,i})=\Big(\frac{d}{2\pi}\Big)^{1/2}e^{-d\eta_{\alpha,i}^2/2},\qquad \sum_i\langle\eta_{\alpha,i}^2\rangle=1,\qquad \langle\eta_{\alpha,i}\eta_{\beta,j}\rangle=\frac1d\delta_{\alpha\beta}\delta_{ij}.\]
高斯的四次平均可以「配對」（Wick）：$\langle\eta_1\eta_2\eta_3\eta_4\rangle=\langle\eta_1\eta_2\rangle\langle\eta_3\eta_4\rangle+\langle\eta_1\eta_3\rangle\langle\eta_2\eta_4\rangle+\langle\eta_1\eta_4\rangle\langle\eta_2\eta_3\rangle$
（一維特例 $\langle x^4\rangle=3\langle x^2\rangle^2$）。$T_{11}^2=\sum a_{\alpha\beta}a_{\gamma\delta}\eta_{\alpha,1}\eta_{\beta,1}\eta_{\gamma,1}\eta_{\delta,1}$ 的三種配對給 $(\mathrm{Tr}\,\mathbf a)^2$、$\mathrm{Tr}\,\mathbf a^2$、$\mathrm{Tr}\,\mathbf a^2$：
\[\langle T_{11}^2\rangle=\frac{(\mathrm{Tr}\,\mathbf a)^2+2\,\mathrm{Tr}\,\mathbf a^2}{d^2},\qquad \langle T_{11}T_{22}\rangle=\frac{(\mathrm{Tr}\,\mathbf a)^2}{d^2},\qquad
\langle T_{12}^2\rangle=\frac{\mathrm{Tr}\,\mathbf a^2}{d^2}\tag{6.33, 6.36, 6.37}\]
（高斯步伐時是精確的）。$N$ 大時 $\mathrm{Tr}\,\mathbf a=\sum_kR_k^2\to\frac N6$、$\mathrm{Tr}\,\mathbf a^2\to\frac{N^2}{90}$（就是 (6.19) 的 $\sum k^{-2}$、$\sum k^{-4}$），所以
\[\langle T_{11}^2\rangle=\frac{N^2}{\pi^4d^2}\Big[\frac{\pi^4}{36}+\frac{2\pi^4}{90}\Big]=\frac{N^2}{20d^2},\qquad \langle T_{11}T_{22}\rangle=\frac{N^2}{36d^2},\qquad \langle T_{12}^2\rangle=\frac{N^2}{90d^2}.\tag{6.35--6.37}\]
比值 $\frac95$、$\frac92$；代入 (6.14)，分子分母乘 180：$\frac{9-5+2d}{9+5(d-1)}=\frac{4+2d}{4+5d}$。\textbf{揭曉}：(6.15) 就這樣算出來了。
\begin{errbox}[p.~141，(6.33a)(6.35a)(6.35b)]
(6.33a) 的和寫成 $\sum_{\alpha,\gamma}a_{\alpha\alpha}a_{\delta\delta}$，應為 $\sum_{\alpha,\delta}$。(6.35a) 多一個 $d^2$，$(\sum_n\frac1{n^2})$ 也少了平方。
(6.35b) 寫 $\frac{\pi^4}{36}+\frac{\pi^4}{90}$，少了 2：照書上得 $\frac{7N^2}{180d^2}$，不是 $\frac{N^2}{20d^2}$。
\end{errbox}

\textbf{對照}（$d=3$）：$N=4$ 的精確 $A=0.6013$，$N=100$ 是 0.5265，$N\to\infty$ 是 $\frac{10}{19}=0.5263$。高斯步伐 2 萬條 $N=100$ 的模擬：
$\langle T_{11}^2\rangle=56.09$、$\langle T_{11}T_{22}\rangle=31.31$、$\langle T_{12}^2\rangle=12.56$（精確 56.67、31.48、12.60），$A=0.5261$。
格子漫步的步伐不是高斯，四次平均多一項（同一步的四次累積量），$N=100$ 的精確值變成 0.5235——就是 Fig.~6.5 那組模擬的對照值。
\textbf{補充 2 的大要}（pp.~159--165）：不用 Kramers 恆等式、直接展開也行：$\langle T_{11}\rangle=\frac{Nl^2}{6d}$、$\langle\mathrm{Tr}\,\mathbf T\rangle=\frac{Nl^2}6$；
$\langle T_{11}^2\rangle$ 的三張圖係數 $2+16+18=36$，$\frac{36}{720}=\frac1{20}$，與上面一致。
\begin{errbox}[p.~159，(6:S2-1)；p.~164]
(6:S2-1) 用比值寫 $A_d$ 時，分母少乘 $\frac{\langle T_{11}T_{22}\rangle}{\langle T_{12}^2\rangle}$。正確是
$A_d=\big[\frac{\langle T_{11}^2\rangle-\langle T_{11}T_{22}\rangle}{\langle T_{12}^2\rangle}+d\big]\big/\big[\big(\frac{\langle T_{11}^2\rangle}{\langle T_{11}T_{22}\rangle}+d-1\big)\frac{\langle T_{11}T_{22}\rangle}{\langle T_{12}^2\rangle}\big]$；
照書上，$d=3$ 得 $\frac{25}{19}$（是 $\frac{10}{19}$ 的 $\frac52$ 倍）。p.~164（(6:S2-25) 前）的「Adding (6:S2-16), (6:S2-21), and (6:S2-23)」應為 (6:S2-24)。
\end{errbox}

\subsubsection{\texorpdfstring{$1/d$}{1/d} 展開：尖峰變寬}
\leavevmode\bk{pp.~142--152}\ \anim{W606}\par\nopagebreak

\textbf{問}：在 100 維，最大的三個主迴轉半徑是固定的數嗎？20 維、3 維呢？

\textbf{配對數}（練習 6.7）：$2m$ 個 $\eta$ 配成 $m$ 對的方法數是 $(2m-1)!!=\frac{(2m)!}{2^mm!}$：$m=1..5$ 是 1、3、15、105、945，正是標準高斯的 $\langle x^{2m}\rangle$。
算 $\langle\mathrm{Tr}\,\mathbf T^n\rangle$ 要把 $2n$ 個 $\eta$ 配對；書上把這些配對按 $\frac1d$ 的次方分類。

\textbf{預解式}：特徵值的分布可以從
\[\mathcal R(\lambda)=\mathrm{Tr}\frac{1}{\lambda\mathbf I-\mathbf T}=\sum_{j=1}^{d}\frac{1}{\lambda-\lambda_j},\qquad
\frac{1}{x-i\epsilon}=P\frac1x+i\pi\delta(x),\qquad \mathrm{Im}\,\mathcal R(\lambda-i\epsilon)=\pi\sum_j\delta(\lambda-\lambda_j)\tag{6.38--6.42}\]
讀出來：虛部在每個特徵值上有一根 $\delta$ 尖峰（$\frac{\epsilon}{x^2+\epsilon^2}$ 底下的面積是 $\pi=3.1416$）。把 $\frac1{\lambda-\mathbf T}=\sum_n\frac{\mathbf T^n}{\lambda^{n+1}}$ 展開、逐項平均，就得到平均的分布。
\begin{errbox}[p.~146，(6.52)]
書上寫 $\mathrm{Im}\,\mathcal R(\lambda-i\epsilon)=\sum_i\delta(\lambda-\lambda_i)+O(1/d)$，漏了 $\pi$（同 (6.42)）：應為 $\pi\sum_i\delta(\lambda-\lambda_i)$。
\end{errbox}
\textbf{零階}（只配相鄰的 $\eta$，(6.50)）：$\langle\mathrm{Tr}(\mathbf T/\lambda)^n\rangle=\mathrm{Tr}(\mathbf a/\lambda)^n$，特徵值就是 $\mathbf a$ 的特徵值——又是樓梯，$\delta$ 尖峰在
$\frac{N+1}{\pi^2k^2}$（除以 $N+1$：0.1013、0.0253、0.0113）。
\textbf{一階}（(6.55)--(6.62)）：尖峰移動、變寬，
\[\langle\Lambda_k\rangle=\frac{N+1}{\pi^2k^2}\Big(1+\frac{3}{4d}\Big),\qquad \big\langle(\Lambda_k-\langle\Lambda_k\rangle)^2\big\rangle=\frac2d\Big(\frac{N+1}{\pi^2k^2}\Big)^2.\tag{6.61, 6.62}\]
位移來自 $\frac1d\sum_{j\ne k}\frac{\lambda_j}{\lambda_k-\lambda_j}=\frac1d\sum_{j\ne k}\frac{k^2}{j^2-k^2}=\frac{3}{4d}$（和 $k$ 無關；$k=1$ 時 $\sum_{j\ge2}\frac1{j^2-1}=\frac12(1+\frac12)=\frac34$）。
相對位移 $\frac3{4d}$：$d=100,20,3$ 是 0.0075、0.0375、0.2500；相對寬度 $\sqrt{2/d}$：0.1414、0.3162、0.8165。

\textbf{模擬}（書上 Fig.~6.12--6.14 的樣子；高斯步伐、$N=50$，$d=100$、20 各 2000 條，$d=3$ 有 2 萬條；圖 \ref{six-fig-oned}）。
\textbf{揭曉}：$d=100$ 的平均與變異數都和 (6.61)(6.62) 吻合到幾個百分點；$d=20$ 的 $k=2,3$ 變異數開始偏小；$d=3$ 只有最大的 $R_1^2$ 還準（平均比 1.01、變異數比 1.01），
$R_2^2$、$R_3^2$ 的變異數只有預測的 0.47、0.23——尖峰寬到互相重疊，排序 $R_1^2\ge R_2^2\ge R_3^2$ 把下面兩個夾住了。
書上 (6.63)--(6.71) 的數字（$d=100$ 變異數比 0.97、0.93、0.95；$d=3$ 是 0.98、0.48、0.23，平均比 1.01、0.91、0.75）和表裡的差都不到 0.05。

\begin{figure}[htbp]
\centering
\SixFigOneDee\\[4pt]
{\footnotesize\SixTabOneDee}
\caption{$1/d$ 展開與模擬（高斯步伐、$N=50$，數值都除以 $N+1$）。上：$\Lambda_k$ 除以零階值的直方圖（藍 $k=1$、橙 $k=2$、綠 $k=3$；縱軸是密度，各圖自己的比例），
灰虛線是零階的 1，黑點線是一階的 $1+\frac3{4d}$：維度越低越寬。下：平均與變異數對 (6.61)(6.62) 的比}\label{six-fig-oned}
\end{figure}

\subsubsection{主迴轉半徑的分布、Fixman、自迴避}
\leavevmode\bk{pp.~152--154}\ \anim{W607}\par\nopagebreak

\textbf{問}：三維漫步最大的 $R_1^2$ 分布長什麼樣子？三個加起來（$R_g^2=\mathrm{Tr}\,\mathbf T$）呢？

\textbf{一個 $R_i^2$}：(6.72) 的積分用最陡下降法估計，得 (6.73)：
\[P(R_i^2)=\frac1{2\pi}\int_{-\infty}^{\infty}\exp\Big[-\frac d2\ln\Big(1+\frac{2ix\alpha_i}{d}\Big)+ixR_i^2\Big]dx,\qquad
P(R_i^2)\propto(R_i^2)^{d/2-1}e^{-dR_i^2/2\alpha_i}.\tag{6.72, 6.73}\]
〔注〕(6.73) 不只是近似：(6.72) 的積分恰好就是這個 Gamma 分布，平均 $\alpha_i$、變異數 $\frac{2\alpha_i^2}d$（練習 6.8）。
$d=3$、$\alpha_1=\frac{N+1}{\pi^2}$（除以 $N+1$ 是 0.1013）：峰在 $\frac{\alpha_1}3$，即 0.0338、高 7.16（圖 \ref{six-fig-fix} 左）。
但模擬的平均是 0.1280——(6.61) 的 0.1267 比 $\alpha_1$ 準得多：三維時峰的位置與寬度都要 $1/d$ 修正。

\textbf{三個加起來：Fixman 的精確結果}（(6.74)--(6.76)）：
\[P(R^2)=\frac1{2\pi}\int_{-\infty}^{\infty}e^{isR^2}K(s)\,ds,\qquad R^2=\sum_{i=1}^dR_i^2,\qquad K(s)=\prod_{n=1}^{N}\Big(1+\frac{2is\alpha_n}{d}\Big)^{-d/2}.\]
為什麼是乘積：$\mathrm{Tr}\,\mathbf T=\sum_n\alpha_n|\vec\xi_n|^2$（$\vec\xi_n$ 是步伐在第 $n$ 個正弦模態上的分量，彼此獨立的高斯），獨立的量相加，特徵函數相乘（練習 6.9）。
用 $\alpha_n=\frac{1}{\pi^2n^2}$（除以 $N+1$）與 $\prod_n(1+\frac{z}{\pi^2n^2})=\frac{\sinh\sqrt z}{\sqrt z}$，得 $K(s)=\big[\sqrt z/\sinh\sqrt z\big]^{d/2}$，$z=\frac{2is}{d}$。
\textbf{算到底}（$d=3$）：平均 $\sum\alpha_n=\frac16=0.1667$、變異數 $\frac2d\sum\alpha_n^2=\frac{1}{45d}=\frac1{135}=0.0074$（標準差 0.0861）。
數值反轉的密度峰在 0.111（高 6.20），$x=0.1,0.2,0.3$ 是 6.078、3.281、1.011；把它積分回去，總和 1.0000、平均 0.1667、變異數 0.0074。
高斯步伐 $N=50$ 的 2 萬條模擬：平均 0.1671、變異數 0.0076（$N=50$ 的精確值 0.1666、0.0074；圖 \ref{six-fig-fix} 右）。
\begin{errbox}[p.~153，(6.76)]
書上寫 $K(R)=\prod_n\big(1-\frac{2is\alpha_n}{d}\big)^{-d/2}$，應為 $K(s)=\prod_n\big(1+\frac{2is\alpha_n}{d}\big)^{-d/2}$（和 (6.72) 同號）。照書上的號，$R^2>0$ 的地方算出來全是 0——密度跑到負的 $R^2$ 去了。
〔注〕要「精確」須用精確的 $\alpha_n$（上面的 $\sin^2$ 公式）；$\alpha_n=\frac{N+1}{\pi^2n^2}$ 是大 $N$ 的極限。
\end{errbox}

\begin{figure}[htbp]
\centering
\SixFigFixman
\caption{三維（$d=3$）主迴轉半徑的分布，高斯步伐 $N=50$ 的 2 萬條模擬，數值除以 $N+1$（30 格、寬 0.02）。
左：最大的 $R_1^2$ 與零階的 Gamma 分布 (6.73)（$\alpha_1=\frac1{\pi^2}$）：寬度差不多，但模擬整個往右移（平均 0.1280 對 $\alpha_1=0.1013$）。右：$R_g^2=\mathrm{Tr}\,\mathbf T$ 與 Fixman 的精確分布 (6.74)：完全吻合}\label{six-fig-fix}
\end{figure}

\textbf{自迴避漫步}（(6.77)，只引用）：不能走回已經走過的點時，$\epsilon=4-d$ 展開的一階結果是
\[A_d=\frac{2d+4}{5d+4}+0.008\,\epsilon,\qquad \epsilon=4-d;\tag{6.77}\]
三維 $0.5263+0.008=0.5343$：自迴避讓漫步稍微更不圓，但影響不大。第 7--12 章會回到自迴避漫步。

% =====================================================================
\subsection{練習 6.1--6.9 解答}
{\small
\begin{enumerate}[label=\textbf{6.\arabic*},itemsep=0.05em]
\item 兩兩差的寫法。\quad 答：記 $n=N+1$ 個點，$\sum_i\sum_j(\vec r_i-\vec r_j)(\vec r_i-\vec r_j)^{\mathsf T}=2n\sum_i\vec r_i\vec r_i^{\mathsf T}-2\big(\sum_i\vec r_i\big)\big(\sum_j\vec r_j\big)^{\mathsf T}$，
除以 $2n^2$ 得 $\frac1n\sum_i\vec r_i\vec r_i^{\mathsf T}-\langle\vec r\rangle\langle\vec r\rangle^{\mathsf T}$，就是 (6.1)。小例子：$\sum_{i<j}(x_i-x_j)^2=14$，除以 $n^2=25$ 正是 $T_{11}=\frac{14}{25}$。
\item 每一步都垂直。\quad 答：$(2d)(2d-2)\cdots(2d-2N+2)=(2d)^N\prod_{k<N}(1-\frac kd)$，$\ln\prod\approx-\sum_{k<N}\frac kd\approx-\frac{N^2}{2d}$；
$d=100$、$N=10$：0.6282 對 0.6065。$N\ll\sqrt d$ 時幾乎全部垂直。
\item 樓梯。\quad 答：(a) $\langle x_n\rangle=\frac{l(N+1-n)}{N+1}$；(b)(c) $n\le m$ 時 $\sum_{j=0}^{N}(x_{j,n}-\langle x_n\rangle)(x_{j,m}-\langle x_m\rangle)=l^2(N+1)\frac{n}{N+1}\big(1-\frac{m}{N+1}\big)$，
除以 $N+1$ 就是 $l^2a_{nm}$；(d) 連續化後積分核 $x_<(1-x_>)$ 的特徵函數 $\sin k\pi y$，特徵值 $\frac{l^2(N+1)}{\pi^2k^2}$（精確 $\frac{l^2}{4(N+1)\sin^2\frac{k\pi}{2(N+1)}}$）；
(e) 所以 $R_k^2=\frac{l^2(N+1)}{\pi^2k^2}$。
\item 反算子。\quad 答：$x<y$ 時 $L=x(1-y)$、$x>y$ 時 $L=y(1-x)$，兩段的二次微分都是 0，在 $x=y$ 斜率從 $1-y$ 跳到 $-y$（跳 $-1$），所以 $-\partial_x^2L=\delta(x-y)$，且 $L(0,y)=L(1,y)=0$。
離散版：$(N+1)\mathbf a$ 是三對角矩陣 $(-1,2,-1)$ 的反矩陣（$N=3$：$\big(\begin{smallmatrix}2&-1&0\\-1&2&-1\\0&-1&2\end{smallmatrix}\big)\cdot\frac14\big(\begin{smallmatrix}3&2&1\\2&4&2\\1&2&3\end{smallmatrix}\big)=\mathbf I$）。
\item 閉合漫步（$\sum_\alpha\vec\eta_\alpha=0$）。\quad 答：把自由步伐扣掉平均步伐就閉合了，共變異變成 $\frac1d(\mathbf I-\frac1N\mathbf J)$（$\mathbf J$ 全是 1）：
$\langle\eta_{\alpha,i}\eta_{\beta,j}\rangle=-\frac{\delta_{ij}}{Nd}$（$\alpha\ne\beta$），對角線 $\frac1d(1-\frac1N)$；每一列加起來是 0，正是 $\sum_\beta\langle\eta_{\alpha,i}\eta_{\beta,j}\rangle=0$。
\item 配對。\quad 答：$\langle T_{11}T_{22}\rangle$ 裡 1 方向和 2 方向的 $\eta$ 不能互配，只剩一種配對：$\frac{(\mathrm{Tr}\,\mathbf a)^2}{d^2}\to\frac{N^2}{36d^2}$；
$\langle T_{12}^2\rangle$ 也只剩一種：$\frac{\mathrm{Tr}\,\mathbf a^2}{d^2}\to\frac{N^2}{90d^2}$。
\item 配對數。\quad 答：第一個 $\eta$ 有 $2m-1$ 個對象可配，剩下的 $2m-2$ 個再配，所以 $(2m-1)(2m-3)\cdots1=\frac{(2m)!}{2^mm!}$：1、3、15、105、945；
高斯的 $\langle x^{2m}\rangle$ 裡每一種配對貢獻 $\langle x^2\rangle^m=1$。
\item Gamma 分布。\quad 答：$P\propto x^{d/2-1}e^{-dx/2\alpha}$ 是形狀 $\frac d2$、尺度 $\frac{2\alpha}d$ 的 Gamma 分布：平均 $\frac d2\cdot\frac{2\alpha}d=\alpha$、變異數 $\frac d2(\frac{2\alpha}d)^2=\frac{2\alpha^2}d$。
\item Fixman。\quad 答：$R^2=\sum_n\alpha_n|\vec\xi_n|^2$ 是獨立的量相加，$\langle e^{-isR^2}\rangle$ 是各自的乘積；每一項是 $\big(1+\frac{2is\alpha_n}d\big)^{-d/2}$（上一題的 Gamma），
所以 $K(s)=\prod_n\big(1+\frac{2is\alpha_n}d\big)^{-d/2}$。
\end{enumerate}}

% =====================================================================
\subsection{第 6 章的回顧與筆誤一覽}
\textbf{回顧}：漫步沒有「標準形狀」，但形狀有統計——三維平均一半扁，維度越高越確定。
\begin{center}
\small\setlength{\tabcolsep}{4pt}
\begin{tabular}{@{}llll@{}}
\toprule
主題 & 結果 & 書上的式子 & 動畫\\
\midrule
形狀 & $\mathbf T$ 的特徵值 $R_i^2$ 與主軸；轉一轉不變 & (6.1)--(6.9) & W601--W602\\
非球度 & $A_d=\frac{4+2d}{4+5d}$，三維 $\frac{10}{19}$；個別值的平均 $A'_3=0.3943$ & (6.13)--(6.17) & W603\\
$d\gg N$ & 樓梯，$R_k^2=\frac{N+1}{\pi^2k^2}$，非球度 $\frac25$ & (6.19)、(6.20) & W604\\
Kramers & $T_{ij}=\sum a_{\alpha\beta}\eta_{\alpha,i}\eta_{\beta,j}$；配對算出 $A_d$ & (6.21)--(6.37) & W605\\
$1/d$ 展開 & 相對位移 $\frac3{4d}$、相對寬度 $\sqrt{2/d}$ & (6.42)、(6.61)、(6.62) & W606\\
分布 & Gamma (6.73)；Fixman 的 $K(s)$；自迴避 0.5343 & (6.72)--(6.77) & W607\\
\bottomrule
\end{tabular}
\end{center}
\textbf{筆誤一覽}（給 2004 年版）：$N$ 與 $N+1$、Fig.~6.4 的線長、(6.19) 的精確 $\sin^2$ 公式、(6.26a) 的和寫到 $N+1$、(6.73) 其實是精確的 Gamma、Fixman 要用精確的 $\alpha_n$ 是〔注〕；
其餘在文中用紅框，每一條都用程式核對過。(6.54) 書上的數字比較像「比值的平均」，不是式子寫的「平均的比」。
\begin{center}
\footnotesize
\SixTabErrata
\end{center}
